\chapter{Variations of Energy}
\label{chap:dc-variations-of-energy}
\section{Formulas for the first and second variations of energy}

\begin{definition}[variation]
  \label{def:dc-ch9-2-1}
  \lean{Riemannian.Variation.IsVariation, Riemannian.Variation.IsVariationOfOrder, Riemannian.Variation.IsSmoothVariation, Riemannian.Variation.IsProperVariation, Riemannian.Variation.IsDifferentiableVariation, Riemannian.Variation.IsDifferentiableVariationOfOrder, Riemannian.Variation.IsSmoothDifferentiableVariation, Riemannian.Variation.curve, Riemannian.Variation.transversal, Riemannian.Variation.variationalField, Riemannian.Variation.variationalFieldBundle, Riemannian.Variation.IsVariationOfOrder.variationalFieldBundle_proj_eq, Riemannian.Variation.IsVariationOfOrder.variationalFieldBundle_piecewise_contMDiff, Riemannian.Variation.IsVariation.curve_zero, Riemannian.Variation.IsVariation.isVariationOfOrderOne, Riemannian.Variation.IsVariation.curve_isPiecewiseDifferentiableCurve, Riemannian.Variation.IsVariationOfOrder.isVariation, Riemannian.Variation.IsVariationOfOrder.curve_piecewise_contMDiff, Riemannian.Variation.IsVariationOfOrder.curve_isPiecewiseDifferentiableCurve, Riemannian.Variation.isProperVariation_iff, Riemannian.Variation.IsDifferentiableVariation.isVariation, Riemannian.Variation.IsDifferentiableVariationOfOrder.isVariationOfOrder}
  \leanok
  \dcref{ch9:2.1}
  Let $c:[0,a]\to M$ be a piecewise differentiable curve in a manifold $M$. A
  \emph{variation} of $c$ is a continuous mapping $f:(-\varepsilon,\varepsilon)\times[0,a]\to M$
  such that:
  \emph{(a)} $f(0,t)=c(t)$, $t\in[0,a]$,
  \emph{(b)} there exists a subdivision of $[0,a]$ by points
  $0=t_0<t_1<\cdots<t_{k+1}=a$, such that the restriction of $f$ to each
  $(-\varepsilon,\varepsilon)\times[t_i,t_{i+1}]$, $i=0,1,\dots,k$, is differentiable.

  A variation is said to be \emph{proper} if $f(s,0)=c(0)$ and $f(s,a)=c(a)$, for all
  $s\in(-\varepsilon,\varepsilon)$. If $f$ is differentiable, the variation is said to be
  \emph{differentiable}. For each $s\in(-\varepsilon,\varepsilon)$, the parametrized curve
  $f_s:[0,a]\to M$ given by $f_s(t)=f(s,t)$ is called a \emph{curve in the variation}, so
  a variation determines a family $f_s(t)$ of neighboring curves of $f_0(t)=c(t)$.
  The variation is proper exactly when the curves in this family have the same
  initial point $c(0)$ and endpoint $c(a)$.
  The parametrized differentiable curve given by $f_t(s)=f(s,t)$, $t$ fixed, is a
  \emph{transversal curve of the variation}. The velocity vector of a transversal curve
  at $s=0$, defined by $V(t)=\frac{\partial f}{\partial s}(0,t)$, is a (piecewise
  differentiable) vector field along $c(t)$ and is called the \emph{variational field} of
  $f$.
\end{definition}
\begin{proposition}[existence of a variation with given variational field]
  \label{prop:dc-ch9-2-2}
  \uses{def:dc-ch9-2-1}
  \lean{Riemannian.Variation.exists_properVariation_with_variationalField}
  \dcref{ch9:2.2}
  Given a piecewise differentiable field $V(t)$, along the piecewise differentiable
  curve $c:[0,a]\to M$, there exists a variation $f:(-\varepsilon,\varepsilon)\times[0,a]\to M$
  of $c$, such that $V(t)$ is the variational field of $f$; in addition, if
  $V(0)=V(a)=0$, it is possible to choose $f$ as a proper variation.
\end{proposition}
\begin{proof}
  \uses{thm:dc-ch3-3-7, lem:dc-ch9-2-2-exponential}
  \sketch
  Since $c([0,a])\subset M$ is compact it is possible to find a $\delta>0$
  such that $\exp_{c(t)}$ is well-defined for all $v\in T_{c(t)}M$ with $|v|<\delta$:
  for each $c(t)$ consider a totally normal neighborhood $W_t$ of $c(t)$ and the
  number $\delta_t>0$ associated to it; a finite subcover
  $W_1,\dots,W_n$ of $c([0,a])$ gives $\delta=\min(\delta_1,\dots,\delta_n)$.
  Let $N=\max_{t\in[0,a]}|V(t)|$, choose $\varepsilon>0$ with
  $\varepsilon N<\delta$, and define
  $f(s,t)=\exp_{c(t)}sV(t)$. Since $\exp_{c(t)}sV(t)=\gamma(1,c(t),sV(t))$ and the
  geodesic depends differentiably on the initial conditions, $f$ is piecewise
  differentiable, and $f(0,t)=c(t)$. The variational field of $f$ is

  $$
  \frac{\partial f}{\partial s}(0,t)=\frac{d}{ds}(\exp_{c(t)}sV(t))\Big|_{s=0}=(d\exp_{c(t)})_0 V(t)=V(t),
  $$

  and if $V(0)=V(a)=0$ then $f$ is proper.
\end{proof}
\begin{definition}[energy of a curve]
  \label{def:dc-ch9-2-2-energy}
  \lean{Riemannian.DCEnergy, Riemannian.dcSpeed, Riemannian.dcArcLength_eq_integral_dcSpeed, Riemannian.dcEnergy_eq_integral_dcSpeed_sq}
  \uses{def:dc-ch1-2-9}
  \leanok
  \dcref{ch9:2.2}
  For a curve $c:[a,b]\to M$ the \emph{speed} is $|dc/dt|=\langle dc/dt,dc/dt\rangle^{1/2}$
  and the \emph{energy} of the segment $c|[a,b]$ is
  $$
  E(c)=\int_a^b\Big|\frac{dc}{dt}\Big|^2\,dt,
  $$
  while its arc length is $L(c)=\int_a^b|dc/dt|\,dt$.
\end{definition}
\begin{lemma}[smooth curves have integrable speed and energy density]
  \label{lem:dc-ch9-2-2-smooth-speed-integrable}
  \lean{Riemannian.Variation.ContMDiffOn.continuousOn_speedSq_of_isOpen, Riemannian.Variation.ContMDiffOn.continuousOn_dcSpeed_of_isOpen, Riemannian.Variation.ContMDiffOn.intervalIntegrable_dcSpeed_of_isOpen, Riemannian.Variation.ContMDiffOn.intervalIntegrable_dcSpeed_sq_of_isOpen}
  \uses{def:dc-ch9-2-2-energy}
  \leanok
  \dcref{ch9:2.2}
  If a curve is smooth on an open neighborhood of $[a,b]$, then its speed and
  squared speed are continuous there and integrable on $[a,b]$. Consequently
  its arc length and energy on $[a,b]$ are represented by ordinary interval
  integrals.
\end{lemma}
\begin{proof}
  \leanok
  \sketch
  Near each time, choose one coordinate chart containing the curve. In that
  chart the velocity depends continuously on time, and the metric coefficients
  are smooth, so the squared speed is continuous. Its nonnegative square root
  is the speed and is also continuous. Continuous functions on the compact
  interval $[a,b]$ are integrable.
\end{proof}
\begin{lemma}[Schwarz inequality $L(c)^2\le a\,E(c)$]
  \label{lem:dc-ch9-2-2-schwarz}
  \lean{Riemannian.dcArcLength_sq_le_mul_dcEnergy, Riemannian.dcArcLength_sq_eq_iff, Riemannian.sq_integral_le_sub_mul_integral_sq, Riemannian.sq_integral_eq_iff_ae_eq_const}
  \uses{def:dc-ch9-2-2-energy}
  \leanok
  \dcref{ch9:2.2}
  For a curve $c:[a,b]\to M$ whose speed and squared speed are integrable,
  $$
  L(c)^2\le (b-a)\,E(c),
  $$
  with equality if and only if the speed $|dc/dt|$ is (almost everywhere) constant, that
  is, if and only if $t$ is proportional to arc length.

  Only integrability is required; the speed of a piecewise differentiable curve
  may jump at subdivision points.
\end{lemma}
\begin{proof}
  \leanok
  \uses{def:dc-ch9-2-2-energy}
  \sketch
  Let $f=|dc/dt|$ and let $m=\frac{1}{b-a}\int_a^b f$ be its mean. Expanding the square,
  $$
  0\le\int_a^b (f-m)^2\,dt=\int_a^b f^2\,dt-2m\int_a^b f\,dt+m^2(b-a)=E(c)-\frac{L(c)^2}{b-a},
  $$
  which is the inequality. Equality holds exactly when $\int_a^b(f-m)^2\,dt=0$, and a
  non-negative integrand has vanishing integral precisely when it vanishes almost
  everywhere; that is, exactly when $f=m$ almost everywhere.
\end{proof}
\begin{lemma}[minimizing geodesics minimize energy]
  \label{lem:dc-ch9-2-3}
  \lean{Riemannian.Geodesic.IsMinimizingGeodesicSegment.dcEnergy_le, Riemannian.Geodesic.IsMinimizingGeodesicSegment.geodesicOn_and_minimizing_of_dcEnergy_eq, Riemannian.Geodesic.IsMinimizingGeodesicSegment.dcEnergy_eq_iff_geodesicOn_and_minimizing}
  \uses{cor:dc-ch3-3-9, lem:dc-ch9-2-2-schwarz, lem:dc-ch9-2-3-inequality, lem:dc-ch9-2-3-equality-case, lem:dc-ch9-2-3-arclength-bridge, lem:dc-ch9-2-3-pointwise-arclength}
  \leanok
  \dcref{ch9:2.3}
  Let $p,q\in M$ and let $\gamma:[0,a]\to M$ be a minimizing geodesic joining $p$ to
  $q$. Then, for all curves $c:[0,a]\to M$ joining $p$ to $q$,

  $$
  E(\gamma)\le E(c)
  $$

  with equality holding if and only if $c$ is a minimizing geodesic.
\end{lemma}
\begin{proof}
  \leanok
  \sketch
  $$
  aE(\gamma)=(L(\gamma))^2\le(L(c))^2\le aE(c),
  $$

  which proves the first assertion. If equality holds, then $(L(c))^2=aE(c)$, so the
  parameter of $c$ is proportional to arc length, and $L(\gamma)=L(c)$, so $c$ is a
  minimizing geodesic (see \cref{cor:dc-ch3-3-9}). The converse is obvious.

\end{proof}
\begin{lemma}[arc length as a path length]
  \label{lem:dc-ch9-2-3-arclength-bridge}
  \lean{Riemannian.ofReal_dcArcLength_eq_pathELength, Riemannian.enorm_mfderiv_eq_ofReal_dcSpeed, Riemannian.dcArcLength_eq_toReal_pathELength, Riemannian.dcArcLength_le_of_pathELength_le}
  \uses{def:dc-ch1-2-9, def:dc-ch9-2-2-energy}
  \leanok
  \dcref{ch9:2.3}
  For a curve $c:[a,b]\to M$ whose speed is integrable, the arc length
  $L(c)=\int_a^b|dc/dt|\,dt$ agrees with the length of $c$
  measured as a path in the metric space $M$. Consequently a comparison of path lengths
  between two curves with integrable speeds is a comparison of their arc lengths $L$.

  (The integrability hypothesis is not removable: a path length is allowed to be infinite,
  while $L(c)=\int_a^b|dc/dt|\,dt$ of a non-integrable speed is not, so for such a curve the
  two genuinely differ.)
\end{lemma}
\begin{proof}
  \leanok
  \uses{def:dc-ch1-2-9}
  \sketch
  The two integrands agree pointwise: the norm of the velocity is $|dc/dt|$ by definition of
  the metric on the tangent spaces. The two notions therefore differ only in the mode of
  integration --- an integral of a real-valued function against $dt$, versus a lower integral
  of the corresponding non-negative extended-real function --- and for a non-negative
  integrable function the two agree.
\end{proof}
\begin{lemma}[a geodesic attains equality in the Schwarz comparison]
  \label{lem:dc-ch9-2-3-geodesic-equality}
  \lean{Riemannian.dcArcLength_sq_eq_mul_dcEnergy_of_isGeodesicOn, Riemannian.dcSpeed_eq_of_isGeodesicOn}
  \uses{lem:dc-ch9-2-2-schwarz, lem:dc-ch3-2-1-constspeed, def:dc-ch9-2-2-energy}
  \leanok
  \dcref{ch9:2.3}
  For a geodesic $\gamma:[a,b]\to M$ with integrable speed and squared speed,
  $$
  L(\gamma)^2 = (b-a)\,E(\gamma).
  $$
\end{lemma}
\begin{proof}
  \leanok
  \uses{lem:dc-ch9-2-2-schwarz, lem:dc-ch3-2-1-constspeed}
  \sketch
  A geodesic has constant speed (\cref{lem:dc-ch3-2-1-constspeed}), so its parameter is
  proportional to arc length, which is precisely the equality case of
  \cref{lem:dc-ch9-2-2-schwarz}.

\end{proof}
\begin{lemma}[a minimizing geodesic minimizes energy: the inequality]
  \label{lem:dc-ch9-2-3-inequality}
  \lean{Riemannian.dcEnergy_le_of_dcArcLength_le, Riemannian.Geodesic.IsMinimizingGeodesicSegment.dcEnergy_le}
  \uses{lem:dc-ch9-2-3-geodesic-equality, lem:dc-ch9-2-2-schwarz, def:dc-ch9-2-2-energy}
  \leanok
  \dcref{ch9:2.3}
  Let $\gamma:[a,b]\to M$ be a geodesic and $c:[a,b]\to M$ a curve with $L(\gamma)\le L(c)$
  (both with integrable speed and squared speed). Then $E(\gamma)\le E(c)$.

\end{lemma}
\begin{proof}
  \leanok
  \uses{lem:dc-ch9-2-3-geodesic-equality, lem:dc-ch9-2-2-schwarz}
  \sketch
  $(b-a)E(\gamma)=L(\gamma)^2\le L(c)^2\le (b-a)E(c)$, the first equality by
  \cref{lem:dc-ch9-2-3-geodesic-equality}, the middle step by hypothesis (both lengths being
  non-negative), and the last by \cref{lem:dc-ch9-2-2-schwarz}. Divide by $b-a>0$.

\end{proof}
\begin{lemma}[a smooth fixed-endpoint variation of a minimizer minimizes energy]
  \label{lem:dc-ch9-2-3-smooth-variation-energy-min}
  \lean{Riemannian.Variation.isLocalMin_dcEnergy_of_smooth_fixedEndpoint_variation}
  \uses{lem:dc-ch9-2-2-smooth-speed-integrable, lem:dc-ch9-2-3-arclength-bridge, lem:dc-ch9-2-3-inequality}
  \leanok
  \dcref{ch9:2.3}
  Let $\gamma:[a,b]\to M$ be a minimizing geodesic, and let
  $f:(-\delta,\delta)\times J\to M$ be a smooth variation on an open
  neighborhood $J$ of $[a,b]$, with $f(0,t)=\gamma(t)$ and with both endpoint
  curves fixed. Then the energy $E(s)=E(f(s,\cdot))$ has a local minimum at
  $s=0$.
\end{lemma}
\begin{proof}
  \leanok
  \sketch
  Every varied slice joins the same endpoints, so its path length is at least
  the distance between them, which is the length of $\gamma$. By
  \cref{lem:dc-ch9-2-2-smooth-speed-integrable}, all speeds and squared speeds
  involved are integrable. The path-length comparison becomes an arc-length
  comparison by \cref{lem:dc-ch9-2-3-arclength-bridge}, and
  \cref{lem:dc-ch9-2-3-inequality} gives $E(0)\le E(s)$ for every sufficiently
  small $s$.
\end{proof}
\begin{lemma}[equality in the energy comparison]
  \label{lem:dc-ch9-2-3-equality-case}
  \lean{Riemannian.dcSpeed_ae_const_of_dcEnergy_eq, Riemannian.dcArcLength_eq_of_dcEnergy_eq}
  \uses{lem:dc-ch9-2-3-geodesic-equality, lem:dc-ch9-2-2-schwarz, def:dc-ch9-2-2-energy}
  \leanok
  \dcref{ch9:2.3}
  In the situation of \cref{lem:dc-ch9-2-3-inequality}, if moreover $E(\gamma)=E(c)$, then
  \emph{(a)} the speed $|dc/dt|$ is (almost everywhere) constant, i.e. the parameter of $c$
  is proportional to arc length; and \emph{(b)} $L(\gamma)=L(c)$, i.e. $c$ is minimizing too.

\end{lemma}
\begin{proof}
  \leanok
  \uses{lem:dc-ch9-2-3-geodesic-equality, lem:dc-ch9-2-2-schwarz}
  \sketch
  Under $E(\gamma)=E(c)$ every inequality in the chain
  $(b-a)E(\gamma)=L(\gamma)^2\le L(c)^2\le (b-a)E(c)$ of
  \cref{lem:dc-ch9-2-3-inequality} collapses to an equality. From $L(c)^2=(b-a)E(c)$ and the
  equality case of \cref{lem:dc-ch9-2-2-schwarz}, the speed of $c$ is a.e. constant, which is
  (a); from $L(\gamma)^2=L(c)^2$ and non-negativity of both lengths, $L(\gamma)=L(c)$, which
  is (b).
\end{proof}
\begin{definition}[the covariant derivative along a curve, chart-independently]
  \label{def:dc-ch9-2-covariant-pair}
  \lean{Riemannian.Variation.IsCovariantDerivSolOn, Riemannian.Variation.IsCovariantDerivFieldAlongOn, Riemannian.Variation.IsCovariantDerivSolOn.transfer, Riemannian.Variation.IsCovariantDerivFieldAlongOn.isCovariantDerivSolOn_of_mem_source}
  \uses{prop:dc-ch2-2-2, lem:dc-ch2-2-2-coord}
  \leanok
  \dcref{ch9:2.4}
  Let $c:I\to M$ be a curve and let $V$, $DV$ be vector fields along $c$. Say that
  \emph{$DV$ is the covariant derivative of $V$ along $c$} on $[a,b]$ if near every
  $t_0\in[a,b]$ there is a chart containing the nearby piece of $c$ in which the
  coordinate readings satisfy formula (1) of \cref{lem:dc-ch2-2-2-coord}, namely
  $\dot V=DV-\Gamma(\dot u,V)(u)$, i.e. $\frac{DV}{dt}=DV$.

  The condition does not depend on the charts chosen: under a change of chart the
  inhomogeneous second-derivative term of the transition map produced by the product rule
  cancels against the transformation law of the Christoffel symbols, which is exactly the
  statement that $\frac{D}{dt}$ is a geometric operator. Being chart-local, the notion is
  meaningful for a curve that leaves every single chart, and on any subinterval lying in
  one chart it localizes back to the coordinate equation.

  (Only differentiability of the curve read in charts is required --- not that $c$ be a
  geodesic. The variations of \cref{def:dc-ch9-2-1} have non-geodesic transversals, so a
  geodesic hypothesis here would exclude precisely the fields this section is about.)
\end{definition}
\begin{lemma}[product rule along a curve, chart-independently]
  \label{lem:dc-ch9-2-covariant-pair-leibniz}
  \lean{Riemannian.Variation.IsCovariantDerivFieldAlongOn.hasDerivAt_metricInner}
  \uses{def:dc-ch9-2-covariant-pair, lem:dc-ch2-3-2-coord}
  \leanok
  \dcref{ch9:2.4}
  Let $V$, $W$ be vector fields along a curve $c$ with covariant derivatives $DV$, $DW$
  in the sense of \cref{def:dc-ch9-2-covariant-pair}. Then, at interior times,
  $$
  \frac{d}{dt}\langle V,W\rangle=\Big\langle\frac{DV}{dt},W\Big\rangle
    +\Big\langle V,\frac{DW}{dt}\Big\rangle.
  $$
\end{lemma}
\begin{proof}
  \leanok
  \uses{def:dc-ch9-2-covariant-pair, lem:dc-ch2-3-2-coord}
  \sketch
  The claim is local in $t$, so fix an interior time and a chart containing the nearby
  piece of $c$, as provided by \cref{def:dc-ch9-2-covariant-pair}. In that chart the
  intrinsic pairing $\langle V,W\rangle$ is the Gram pairing of the coordinate readings,
  and the coordinate metric compatibility of $\frac{D}{dt}$ (\cref{lem:dc-ch2-3-2-coord},
  the Levi-Civita case of formula (3) of \cref{prop:dc-ch2-3-2}) gives its derivative as
  $\langle\frac{DV}{dt},W\rangle+\langle V,\frac{DW}{dt}\rangle$ with $\frac{D}{dt}$ the
  coordinate covariant derivative of \cref{lem:dc-ch2-2-2-coord}. By
  \cref{def:dc-ch9-2-covariant-pair} those two coordinate covariant derivatives are the
  readings of $DV$ and $DW$, which is the assertion.
\end{proof}
\begin{lemma}[the symmetry lemma, operator form]
  \label{lem:dc-ch9-2-4-symmetry}
  \lean{Riemannian.Variation.surfaceCovariantDerivS_snd_eq_surfaceCovariantDerivT_fst, Riemannian.Variation.surfaceCovariantDerivS_snd_eq_surfaceCovariantDerivT_fst_of_forall}
  \uses{lem:dc-ch3-3-4, def:dc-ch3-3-3}
  \leanok
  \dcref{ch9:2.4}
  Let $f$ be a ($C^2$) parametrized surface (\cref{def:dc-ch3-3-3}) \emph{read in a fixed
  chart}, and let $\frac{D}{\partial s}$, $\frac{D}{\partial t}$ be the covariant
  derivatives along $f$ of that chart. Then the two exchange on the velocity fields of $f$:
  $$
  \frac{D}{\partial s}\frac{\partial f}{\partial t}
    = \frac{D}{\partial t}\frac{\partial f}{\partial s}.
  $$
  No curvature term appears. Contrast \cref{lem:dc-ch4-4-1}, which exchanges
  $\frac{D}{\partial s}$ and $\frac{D}{\partial t}$ on an \emph{arbitrary} field $V$ along
  $f$ and does pick up $R(\frac{\partial f}{\partial s},\frac{\partial f}{\partial t})V$;
  here the two sides differ only by the mixed second partial, which Schwarz kills.

\end{lemma}
\begin{proof}
  \leanok
  \uses{lem:dc-ch3-3-4}
  \sketch
  Both sides expand, in a chart, as a mixed second partial of $f$ plus a Christoffel term:
  $\frac{D}{\partial s}\frac{\partial f}{\partial t}
   = \frac{\partial^2 f}{\partial s\,\partial t}
     + \Gamma\big(\frac{\partial f}{\partial s},\frac{\partial f}{\partial t}\big)(f)$, and
  symmetrically with $s$, $t$ interchanged. The mixed partials agree by Schwarz, and the two
  Christoffel terms agree because $\Gamma$ is symmetric in its two vector slots --- which is
  exactly the hypothesis that the connection is symmetric. This is \cref{lem:dc-ch3-3-4} at
  $v=(1,0)$, $w=(0,1)$.
\end{proof}
\begin{lemma}[the $s$-derivative of the energy density]
  \label{lem:dc-ch9-2-4-energy-density}
  \lean{Riemannian.Variation.hasDerivAt_chartEnergyDensity_slice, Riemannian.Variation.hasDerivAt_chartEnergyDensity_slice_symm}
  \uses{lem:dc-ch9-2-4-symmetry, lem:dc-ch2-3-2-coord}
  \leanok
  \dcref{ch9:2.4}
  Let $f$ be a ($C^2$) parametrized surface \emph{read in a fixed chart}, as in
  \cref{lem:dc-ch9-2-4-symmetry}. Then, pointwise,
  $$
  \frac{\partial}{\partial s}\Big\langle\frac{\partial f}{\partial t},\frac{\partial f}{\partial t}\Big\rangle
    = 2\Big\langle\frac{D}{\partial s}\frac{\partial f}{\partial t},\frac{\partial f}{\partial t}\Big\rangle
    = 2\Big\langle\frac{D}{\partial t}\frac{\partial f}{\partial s},\frac{\partial f}{\partial t}\Big\rangle.
  $$
\end{lemma}
\begin{proof}
  \leanok
  \uses{lem:dc-ch9-2-4-symmetry, lem:dc-ch2-3-2-coord}
  \sketch
  The first equality is coordinate metric compatibility (\cref{lem:dc-ch2-3-2-coord})
  applied along the $s$-slice with $V=W=\frac{\partial f}{\partial t}$: the product rule's
  two terms $\langle\frac{D}{\partial s}V,W\rangle+\langle V,\frac{D}{\partial s}W\rangle$
  coincide because the pairing is symmetric, which supplies the factor $2$. The second
  equality is \cref{lem:dc-ch9-2-4-symmetry}, the symmetry of the connection, applied to
  the first factor. No curvature appears.
\end{proof}
\begin{lemma}[integration by parts along a curve]
  \label{lem:dc-ch9-2-4-parts}
  \lean{Riemannian.Variation.IsCovariantDerivFieldAlongOn.integral_metricInner_add, Riemannian.Variation.IsCovariantDerivFieldAlongOn.integral_metricInner_covariantDeriv_left, Riemannian.Variation.IsCovariantDerivFieldAlongOn.integral_metricInner_eq_neg_integral_of_proper}
  \uses{lem:dc-ch9-2-covariant-pair-leibniz, def:dc-ch9-2-covariant-pair}
  \leanok
  \dcref{ch9:2.4}
  Let $V$, $W$ be vector fields along a curve $c$ with covariant derivatives $DV$, $DW$
  in the sense of \cref{def:dc-ch9-2-covariant-pair}, on a segment $[a,b]$ carrying no
  breakpoint, and suppose the pairings $\langle\frac{DV}{dt},W\rangle$ and
  $\langle V,\frac{DW}{dt}\rangle$ are integrable on $[a,b]$. Then
  $$
  \int_a^b\Big\langle\frac{DV}{dt},W\Big\rangle dt
    = \langle V(b),W(b)\rangle-\langle V(a),W(a)\rangle
      -\int_a^b\Big\langle V,\frac{DW}{dt}\Big\rangle dt.
  $$
  If $V(a)=V(b)=0$, the identity becomes
  $\int_a^b\langle\frac{DV}{dt},W\rangle dt=-\int_a^b\langle V,\frac{DW}{dt}\rangle dt$.

  For $W=dc/dt$, summing over the segments of a piecewise differentiable curve
  telescopes the endpoint contributions into
  $-\sum_i\langle V(t_i),\frac{dc}{dt}(t_i^+)-\frac{dc}{dt}(t_i^-)\rangle$.
\end{lemma}
\begin{proof}
  \leanok
  \uses{lem:dc-ch9-2-covariant-pair-leibniz}
  \sketch
  By \cref{lem:dc-ch9-2-covariant-pair-leibniz}, $\frac{d}{dt}\langle V,W\rangle
  =\langle\frac{DV}{dt},W\rangle+\langle V,\frac{DW}{dt}\rangle$ at interior times, and
  $t\mapsto\langle V,W\rangle$ is continuous on $[a,b]$. The fundamental theorem of
  calculus (in the form requiring continuity on $[a,b]$ and a derivative on $(a,b)$ only,
  which is what \cref{lem:dc-ch9-2-covariant-pair-leibniz} supplies) integrates this to
  $\int_a^b(\langle\frac{DV}{dt},W\rangle+\langle V,\frac{DW}{dt}\rangle)dt
  =\langle V(b),W(b)\rangle-\langle V(a),W(a)\rangle$; splitting the integral and
  rearranging gives the claim. When $V(a)=V(b)=0$ both boundary terms vanish by linearity
  of the pairing in its first slot.
\end{proof}
\begin{lemma}[the $s$-derivative of the energy density, chart-free]
  \label{lem:dc-ch9-2-4-density-self}
  \lean{Riemannian.Variation.IsCovariantDerivFieldAlongOn.hasDerivAt_metricInner_self}
  \uses{lem:dc-ch9-2-covariant-pair-leibniz, def:dc-ch9-2-covariant-pair}
  \leanok
  \dcref{ch9:2.4}
  Let $V$ be a vector field along a curve $\gamma$ with covariant derivative $DV$ in the
  sense of \cref{def:dc-ch9-2-covariant-pair}. Then, at interior times,
  $$
  \frac{d}{dt}\langle V,V\rangle = 2\Big\langle\frac{DV}{dt},V\Big\rangle.
  $$

  Applied along a \emph{transversal} $\sigma\mapsto f(\sigma,t)$ of a variation
  (\cref{def:dc-ch9-2-1}) with $V=\frac{\partial f}{\partial t}$, this is the first
  equality of \cref{lem:dc-ch9-2-4-energy-density} with the chart removed: a transversal is
  a curve like any other, and \cref{def:dc-ch9-2-covariant-pair} is chart-independent.
\end{lemma}
\begin{proof}
  \leanok
  \uses{lem:dc-ch9-2-covariant-pair-leibniz}
  \sketch
  By \cref{lem:dc-ch9-2-covariant-pair-leibniz} at $W=V$, one has
  $\frac{d}{dt}\langle V,V\rangle
  =\langle\frac{DV}{dt},V\rangle+\langle V,\frac{DV}{dt}\rangle$ at interior times. The two
  summands are equal by the symmetry of the metric, which supplies the factor $2$.

\end{proof}
\begin{lemma}[symmetry of the connection along a surface, chart-free]
  \label{lem:dc-ch9-2-4-symmetry-manifold}
  \lean{Riemannian.Variation.covariantDerivS_velT_eq_covariantDerivT_velS}
  \uses{lem:dc-ch9-2-4-symmetry, def:dc-ch9-2-covariant-pair}
  \leanok
  \dcref{ch9:2.4}
  Let $f$ be a parametrized surface in $M$, let $\frac{\partial f}{\partial t}$ and
  $\frac{\partial f}{\partial s}$ be its coordinate velocity fields, and suppose
  $\big(\frac{\partial f}{\partial t},\frac{D}{\partial s}\frac{\partial f}{\partial t}\big)$
  is a covariant pair along the transversal through $(s_0,t_0)$ and
  $\big(\frac{\partial f}{\partial s},\frac{D}{\partial t}\frac{\partial f}{\partial s}\big)$
  a covariant pair along the curve in the variation through $(s_0,t_0)$, both in the sense
  of \cref{def:dc-ch9-2-covariant-pair}. If $f$ is $C^2$ near $(s_0,t_0)$ and the two slice
  curves through $(s_0,t_0)$ stay in the source of a single chart on their windows, then
  $$
  \frac{D}{\partial s}\frac{\partial f}{\partial t}(s_0,t_0)
    = \frac{D}{\partial t}\frac{\partial f}{\partial s}(s_0,t_0).
  $$

  The hypotheses concern only the two \emph{slice lines} through $(s_0,t_0)$.
\end{lemma}
\begin{proof}
  \leanok
  \uses{lem:dc-ch9-2-4-symmetry, def:dc-ch9-2-covariant-pair}
  \sketch
  Choose a chart containing the relevant pieces of both slice curves. In that chart,
  the covariant-pair hypotheses identify the intrinsic derivatives with the coordinate
  covariant derivatives. Apply \cref{lem:dc-ch9-2-4-symmetry} to the coordinate
  representation and transport the resulting equality back to $T_{f(s_0,t_0)}M$.

\end{proof}
\begin{lemma}[differentiation of the energy under the integral sign]
  \label{lem:dc-ch9-2-4-under-integral}
  \lean{Riemannian.Variation.hasDerivAt_dcEnergy_of_dominated}
  \uses{lem:dc-ch9-2-4-density-self, def:dc-ch9-2-2-energy, def:dc-ch9-2-1, def:dc-ch9-2-covariant-pair}
  \leanok
  \dcref{ch9:2.4}
  Let $f$ be a variation of $c$ and let $E$ be its energy (\cref{def:dc-ch9-2-2-energy}).
  Suppose $\big(\frac{\partial f}{\partial t},\frac{D}{\partial s}\frac{\partial f}{\partial t}\big)$
  is a covariant pair along each transversal, that the energy density is integrable at
  $s_0$, and that on a neighbourhood of $s_0$ the integrands
  $\big\langle\frac{D}{\partial s}\frac{\partial f}{\partial t},\frac{\partial f}{\partial t}\big\rangle$
  are dominated by a fixed integrable function of $t$. Then
  $$
  E'(s_0)
    = 2\int_a^b\Big\langle\frac{D}{\partial s}\frac{\partial f}{\partial t},\frac{\partial f}{\partial t}\Big\rangle\Big|_{s_0} dt .
  $$

\end{lemma}
\begin{proof}
  \leanok
  \uses{lem:dc-ch9-2-4-density-self, def:dc-ch9-2-2-energy}
  \sketch
  By \cref{def:dc-ch9-2-2-energy}, $E(s)=\int_a^b\langle\frac{\partial f}{\partial t},\frac{\partial f}{\partial t}\rangle dt$
  with the integrand evaluated along the curve $f(s,\cdot)$. For each fixed $t$ the map
  $\sigma\mapsto\langle\frac{\partial f}{\partial t},\frac{\partial f}{\partial t}\rangle(\sigma,t)$
  is, by \cref{lem:dc-ch9-2-4-density-self} applied along the transversal through $t$,
  differentiable at each interior $\sigma$ with derivative
  $2\langle\frac{D}{\partial s}\frac{\partial f}{\partial t},\frac{\partial f}{\partial t}\rangle(\sigma,t)$.
  Together with integrability at $s_0$ and the domination hypothesis, the standard theorem
  on differentiation under the integral sign applies and yields the displayed identity.

\end{proof}
\begin{lemma}[first variation on a segment without breakpoints]
  \label{lem:dc-ch9-2-4-segment}
  \lean{Riemannian.Variation.hasDerivAt_dcEnergy_eq_first_variation}
  \uses{lem:dc-ch9-2-4-under-integral, lem:dc-ch9-2-4-parts, def:dc-ch9-2-2-energy, def:dc-ch9-2-covariant-pair}
  \leanok
  \dcref{ch9:2.4}
  In the situation of \cref{lem:dc-ch9-2-4-under-integral}, let $V=\frac{\partial f}{\partial s}(s_0,\cdot)$
  carry the covariant derivative $\frac{DV}{dt}$ and let $\frac{dc}{dt}=\frac{\partial f}{\partial t}(s_0,\cdot)$
  carry $\frac{D}{dt}\frac{dc}{dt}$, both along $c=f(s_0,\cdot)$ and in the sense of
  \cref{def:dc-ch9-2-covariant-pair}. Assume the symmetry of the connection along $c$,
  $$
  \frac{D}{\partial s}\frac{\partial f}{\partial t}(s_0,t)=\frac{D}{\partial t}\frac{\partial f}{\partial s}(s_0,t)
  \qquad\text{for every interior }t\in(a,b),
  $$
  and the integrability of the two pairings of \cref{lem:dc-ch9-2-4-parts}. Then
  $$
  \frac{1}{2}E'(s_0)
    = -\int_a^b\Big\langle V,\frac{D}{dt}\frac{dc}{dt}\Big\rangle dt
      -\Big\langle V(a),\frac{dc}{dt}(a)\Big\rangle
      +\Big\langle V(b),\frac{dc}{dt}(b)\Big\rangle .
  $$

\end{lemma}
\begin{proof}
  \leanok
  \uses{lem:dc-ch9-2-4-under-integral, lem:dc-ch9-2-4-parts}
  \sketch
  By \cref{lem:dc-ch9-2-4-under-integral},
  $E'(s_0)=2\int_a^b\langle\frac{D}{\partial s}\frac{\partial f}{\partial t},\frac{\partial f}{\partial t}\rangle dt$.
  The symmetry hypothesis rewrites the integrand as
  $\langle\frac{DV}{dt},\frac{dc}{dt}\rangle$ pointwise, so
  $\frac{1}{2}E'(s_0)=\int_a^b\langle\frac{DV}{dt},\frac{dc}{dt}\rangle dt$. Applying
  \cref{lem:dc-ch9-2-4-parts} with $W=\frac{dc}{dt}$ and $DW=\frac{D}{dt}\frac{dc}{dt}$
  turns the right-hand side into
  $\langle V(b),\frac{dc}{dt}(b)\rangle-\langle V(a),\frac{dc}{dt}(a)\rangle
  -\int_a^b\langle V,\frac{D}{dt}\frac{dc}{dt}\rangle dt$, which is the claim.
\end{proof}
\begin{proposition}[formula for the first variation of the energy of a curve]
  \label{prop:dc-ch9-2-4}
  \uses{def:dc-ch9-2-1, def:dc-ch9-2-2-energy, lem:dc-ch9-2-4-symmetry, def:dc-ch9-2-covariant-pair, lem:dc-ch9-2-4-energy-density, lem:dc-ch9-2-4-parts, lem:dc-ch9-2-4-density-self, lem:dc-ch9-2-4-symmetry-manifold, lem:dc-ch9-2-4-under-integral, lem:dc-ch9-2-4-segment, lem:dc-ch9-2-4-piecewise-assembly}
  \dcref{ch9:2.4}
  Let $c:[0,a]\to M$ be a piecewise differentiable curve and let
  $f:(-\varepsilon,\varepsilon)\times[0,a]\to M$ be a variation of $c$. If
  $E:(-\varepsilon,\varepsilon)\to\mathbf{R}$ is the energy of $f$ then

  $$
  \frac{1}{2}E'(0)=-\int_0^a\Big\langle V(t),\frac{D}{dt}\frac{dc}{dt}\Big\rangle dt-\sum_{i=1}^k\Big\langle V(t_i),\frac{dc}{dt}(t_i^+)-\frac{dc}{dt}(t_i^-)\Big\rangle-\Big\langle V(0),\frac{dc}{dt}(0)\Big\rangle+\Big\langle V(a),\frac{dc}{dt}(a)\Big\rangle,\qquad(1)
  $$

  where $V(t)$ is the variational field of $f$, and
  $\frac{dc}{dt}(t_i^+)=\lim_{t\to t_i,\,t>t_i}\frac{dc}{dt}$,
  $\frac{dc}{dt}(t_i^-)=\lim_{t\to t_i,\,t<t_i}\frac{dc}{dt}$.
\end{proposition}
\begin{proof}
  \uses{def:dc-ch9-2-2-energy, lem:dc-ch9-2-4-energy-density, lem:dc-ch9-2-4-parts,
    def:dc-ch9-2-covariant-pair}
  \sketch
  By definition,
  $E(s)=\int_0^a\langle\frac{\partial f}{\partial t},\frac{\partial f}{\partial t}\rangle dt=\sum_{i=0}^k\int_{t_i}^{t_{i+1}}\langle\frac{\partial f}{\partial t},\frac{\partial f}{\partial t}\rangle dt$.
  Differentiating under the integral sign and using the symmetry of the Riemannian
  connection,

  $$
  \frac{d}{ds}\int_{t_i}^{t_{i+1}}\Big\langle\frac{\partial f}{\partial t},\frac{\partial f}{\partial t}\Big\rangle dt=2\Big\langle\frac{\partial f}{\partial s},\frac{\partial f}{\partial t}\Big\rangle\Big|_{t_i}^{t_{i+1}}-2\int_{t_i}^{t_{i+1}}\Big\langle\frac{\partial f}{\partial s},\frac{D}{dt}\frac{\partial f}{\partial t}\Big\rangle dt.
  $$

  Therefore,

  $$
  \frac{1}{2}\frac{dE}{ds}=\sum_{i=0}^k\Big\langle\frac{\partial f}{\partial s},\frac{\partial f}{\partial t}\Big\rangle\Big|_{t_i}^{t_{i+1}}-\int_0^a\Big\langle\frac{\partial f}{\partial s},\frac{D}{dt}\frac{\partial f}{\partial t}\Big\rangle dt.\qquad(2)
  $$

  Putting $s=0$ in (2) yields (1).
\end{proof}
\begin{lemma}[a geodesic is a critical point of the energy]
  \label{lem:dc-ch9-2-5-geodesic}
  \lean{Riemannian.Variation.IsCovariantDerivFieldAlongOn.integral_metricInner_covariantDeriv_eq_zero_of_geodesic}
  \uses{lem:dc-ch9-2-4-parts, def:dc-ch9-2-covariant-pair}
  \leanok
  \dcref{ch9:2.5}
  Let $W$ be a \emph{parallel} field along a curve $c$ --- $\frac{DW}{dt}=0$ --- and let
  $V$, $DV$ be a covariant-derivative pair along $c$ with $V(a)=V(b)=0$. Then on the
  segment $[a,b]$
  $$
  \int_a^b\Big\langle\frac{DV}{dt},W\Big\rangle dt=0.
  $$

  For $W=dc/dt$, parallelity is the geodesic equation; for the variational
  field of a proper variation, the endpoint condition holds. Together with
  \cref{prop:dc-ch9-2-4}, this gives the forward implication of
  \cref{prop:dc-ch9-2-5}.
\end{lemma}
\begin{proof}
  \leanok
  \uses{lem:dc-ch9-2-4-parts}
  \sketch
  Apply \cref{lem:dc-ch9-2-4-parts} with $DW=0$. The remaining integral
  $\int_a^b\langle V,\frac{DW}{dt}\rangle dt$ vanishes because its integrand is
  identically zero, and the two boundary terms vanish because $V(a)=V(b)=0$.
\end{proof}
\begin{proposition}[geodesics as critical points of energy]
  \label{prop:dc-ch9-2-5}
  \uses{prop:dc-ch9-2-4, prop:dc-ch9-2-2, lem:dc-ch9-2-5-geodesic}
  \dcref{ch9:2.5}
  A piecewise differentiable curve $c:[0,a]\to M$ is a geodesic if and only if, for
  every proper variation $f$ of $c$, we have $\frac{dE}{ds}(0)=0$.
\end{proposition}
\begin{proof}
  \uses{prop:dc-ch9-2-4, prop:dc-ch9-2-2, lem:dc-ch9-2-5-geodesic}
  \sketch
  If $c$ is a geodesic, $\frac{D}{dt}\frac{dc}{dt}=0$ and $c$ is regular;
  for $f$ proper, $V(0)=V(a)=0$, and all terms of (1) are zero
  (\cref{lem:dc-ch9-2-5-geodesic}). Conversely, suppose
  $\frac{dE}{ds}(0)=0$ for all proper variations. Let
  $V(t)=g(t)\frac{D}{dt}\frac{dc}{dt}$, where $g:[0,a]\to\mathbf{R}$ is piecewise
  differentiable with $g(t)>0$ if $t\ne t_i$ and $g(t_i)=0$, $i=0,\dots,k+1$.
  Constructing a variation
  with this $V(t)$,

  $$
  \frac{1}{2}\frac{dE}{ds}(0)=-\int_0^a g(t)\Big\langle\frac{D}{dt}\frac{dc}{dt},\frac{D}{dt}\frac{dc}{dt}\Big\rangle dt=0,
  $$

  so $\frac{D}{dt}\frac{dc}{dt}=0$ on each $(t_i,t_{i+1})$, that is, $c$ is a geodesic
  on each $(t_i,t_{i+1})$. To handle the points $t_i$, take another variational field
  $\overline V(t)$ with $\overline V(0)=\overline V(a)=0$ and, for $t\ne 0,a$,
  $\overline V(t_i)=\frac{dc}{dt}(t_i^+)-\frac{dc}{dt}(t_i^-)$. Then

  $$
  0=\frac{1}{2}\frac{dE}{ds}(0)=-\sum_{i=1}^k\Big|\frac{dc}{dt}(t_i^+)-\frac{dc}{dt}(t_i^-)\Big|^2,
  $$

  so $c$ is of class $C^1$ at each $t_i$. Since $\frac{D}{dt}\frac{dc}{dt}=0$ at
  $t_i$, $c$ satisfies the geodesic equation on $(0,a)$; by uniqueness of solutions
  to ordinary differential equations, $c\in C^\infty$ and is a geodesic.
\end{proof}
\begin{remark}[Path-space interpretation of the energy functional]
  \label{rem:dc-ch9-2-7}
  \dcref{ch9:2.7}
  For fixed $p,q\in M$, let $\Omega_{p,q}$ denote the space of piecewise
  differentiable curves from $p$ to $q$. Formally, its tangent space at a curve
  $c$ consists of piecewise differentiable vector fields along $c$ that vanish
  at both endpoints. The first variation is the directional derivative of
  $E:\Omega_{p,q}\to\mathbf{R}$, and the geodesics from $p$ to $q$ are its
  critical points. This path space is infinite-dimensional.
\end{remark}
\begin{lemma}[the second differentiation under the integral sign]
  \label{lem:dc-ch9-2-8-under-integral}
  \lean{Riemannian.Variation.hasDerivAt_dcPairing_of_dominated}
  \uses{lem:dc-ch9-2-covariant-pair-leibniz, def:dc-ch9-2-covariant-pair, lem:dc-ch9-2-4-under-integral}
  \leanok
  \dcref{ch9:2.8}
  Let $f$ be a variation of $c$, let $\frac{\partial f}{\partial t}$ carry the covariant
  $s$-derivative $\frac{D}{\partial s}\frac{\partial f}{\partial t}$ along each transversal,
  and let $\frac{D}{\partial s}\frac{\partial f}{\partial t}$ in turn carry its own covariant
  $s$-derivative $\frac{D}{\partial s}\frac{D}{\partial s}\frac{\partial f}{\partial t}$, all
  in the sense of \cref{def:dc-ch9-2-covariant-pair}; suppose the second integrand is
  dominated near $s_0$ by a fixed integrable function of $t$. Then
  $$
  \frac{d}{ds}\int_a^b\Big\langle\frac{D}{\partial s}\frac{\partial f}{\partial t},
    \frac{\partial f}{\partial t}\Big\rangle dt\Big|_{s_0}
  = \int_a^b\Big\{\Big\langle\frac{D}{\partial s}\frac{D}{\partial s}\frac{\partial f}{\partial t},
    \frac{\partial f}{\partial t}\Big\rangle
    + \Big\langle\frac{D}{\partial s}\frac{\partial f}{\partial t},
      \frac{D}{\partial s}\frac{\partial f}{\partial t}\Big\rangle\Big\}\,dt .
  $$
  The pointwise input is metric compatibility along each transversal
  (\cref{lem:dc-ch9-2-covariant-pair-leibniz}) applied to the pairs
  $\big(\frac{D}{\partial s}\frac{\partial f}{\partial t},
    \frac{D}{\partial s}\frac{D}{\partial s}\frac{\partial f}{\partial t}\big)$ and
  $\big(\frac{\partial f}{\partial t},\frac{D}{\partial s}\frac{\partial f}{\partial t}\big)$;
  the exchange of $\frac{d}{ds}$ with $\int dt$ is the same differentiation-under-the-integral
  theorem used in \cref{lem:dc-ch9-2-4-under-integral}.
\end{lemma}
\begin{proof}
  \leanok
  \uses{lem:dc-ch9-2-covariant-pair-leibniz}
  \sketch
  For each fixed $t$ the map
  $\sigma\mapsto\langle\frac{D}{\partial s}\frac{\partial f}{\partial t},\frac{\partial f}{\partial t}\rangle(\sigma,t)$
  is differentiable at interior $\sigma$ with derivative
  $\langle\frac{D}{\partial s}\frac{D}{\partial s}\frac{\partial f}{\partial t},\frac{\partial f}{\partial t}\rangle
  +\langle\frac{D}{\partial s}\frac{\partial f}{\partial t},\frac{D}{\partial s}\frac{\partial f}{\partial t}\rangle$,
  by \cref{lem:dc-ch9-2-covariant-pair-leibniz} (the product rule, not its $V=W$ collapse).
  Together with the domination hypothesis the standard theorem on differentiation under the
  integral sign yields the displayed identity.
\end{proof}
\begin{lemma}[the second variation of the energy, before the curvature substitution]
  \label{lem:dc-ch9-2-8-energy-snd-deriv}
  \lean{Riemannian.Variation.hasDerivAt_deriv_dcEnergy_second_variation}
  \uses{lem:dc-ch9-2-8-under-integral, lem:dc-ch9-2-4-under-integral, def:dc-ch9-2-2-energy, def:dc-ch9-2-covariant-pair}
  \leanok
  \dcref{ch9:2.8}
  In the situation of \cref{lem:dc-ch9-2-8-under-integral}, suppose moreover that the first
  variation $E'(\sigma)=2\int_a^b\langle\frac{D}{\partial s}\frac{\partial f}{\partial t},
  \frac{\partial f}{\partial t}\rangle\,dt$ of \cref{lem:dc-ch9-2-4-under-integral} holds at
  every $\sigma$ in a neighbourhood of $s_0$. Then
  $$
  \tfrac12 E''(s_0)
  = \int_a^b\Big\{\Big\langle\frac{D}{\partial s}\frac{D}{\partial s}\frac{\partial f}{\partial t},
    \frac{\partial f}{\partial t}\Big\rangle
    + \Big\langle\frac{D}{\partial s}\frac{\partial f}{\partial t},
      \frac{D}{\partial s}\frac{\partial f}{\partial t}\Big\rangle\Big\}\,dt .
  $$

\end{lemma}
\begin{proof}
  \leanok
  \uses{lem:dc-ch9-2-8-under-integral, lem:dc-ch9-2-4-under-integral}
  \sketch
  The first variation, supplied on a neighbourhood of $s_0$, says that $E'$ agrees there with
  $\sigma\mapsto 2\int_a^b\langle\frac{D}{\partial s}\frac{\partial f}{\partial t},\frac{\partial f}{\partial t}\rangle\,dt$;
  hence $E''(s_0)$ is the derivative of the latter at $s_0$, which
  \cref{lem:dc-ch9-2-8-under-integral} computes. Transferring the derivative along the
  eventual equality gives the claim.
\end{proof}
\begin{lemma}[the Ricci-identity substitution, metric-paired, chart-read]
  \label{lem:dc-ch9-2-8-curvature-substitution}
  \lean{Riemannian.Variation.chartMetricInner_surface_covariant_commutator_eq_curvatureFormAt, Riemannian.Variation.chartMetricInner_chartCurvatureContraction2_eq_curvatureFormAt, Riemannian.Variation.chartMetricInner_chartCurvatureContraction2_chartVectorRep_eq_curvatureFormAt}
  \uses{lem:dc-ch4-4-1, def:dc-ch4-2-1}
  \leanok
  \dcref{ch9:2.8}
  Let $f$ be a ($C^2$) parametrized surface \emph{read in a fixed chart} and $V$ a field along
  it. The commutator of the two covariant derivatives, paired against a vector $w$, is the
  intrinsic curvature $(0,4)$ form of the chart-frame realizations of $\frac{\partial f}{\partial s}$,
  $\frac{\partial f}{\partial t}$, $V$ and $w$:
  $$
  \Big\langle\frac{D}{\partial t}\frac{D}{\partial s}V - \frac{D}{\partial s}\frac{D}{\partial t}V,\ w\Big\rangle
    = \Big\langle R\Big(\frac{\partial f}{\partial s},\frac{\partial f}{\partial t}\Big)V,\ w\Big\rangle,
  $$
\end{lemma}
\begin{proof}
  \leanok
  \uses{lem:dc-ch4-4-1}
  \sketch
  The chart-level Ricci identity (\cref{lem:dc-ch4-4-1}) equates the commutator with the
  coordinate curvature contraction $\mathcal{R}_{\mathrm{chart}}(\frac{\partial f}{\partial s},\frac{\partial f}{\partial t})V$.
  Pairing against $w$ and using the coordinate-to-intrinsic curvature identity rewrites
  the result as the intrinsic curvature $(0,4)$ form
  $R(\frac{\partial f}{\partial s},\frac{\partial f}{\partial t},V,w)$.
\end{proof}
\begin{lemma}[the Ricci-identity substitution, chart-free]
  \label{lem:dc-ch9-2-8-curvature-substitution-manifold}
  \lean{Riemannian.Variation.metricInner_covariantDerivT_covariantDerivS_commutator_eq_curvatureFormAt}
  \uses{lem:dc-ch9-2-8-curvature-substitution, def:dc-ch9-2-covariant-pair, lem:dc-ch9-2-4-symmetry-manifold}
  \leanok
  \dcref{ch9:2.8}
  Let $f$ be a parametrized surface in $M$ and $V$ a field along it, and suppose
  $\frac{D}{\partial s}V$, $\frac{D}{\partial t}V$ are covariant derivatives of $V$ along the
  slices in the sense of \cref{def:dc-ch9-2-covariant-pair}, and
  $\frac{D}{\partial t}\frac{D}{\partial s}V$, $\frac{D}{\partial s}\frac{D}{\partial t}V$ the
  covariant derivatives of those, along the two slice curves through $(s_0,t_0)$. Then, at
  $(s_0,t_0)$,
  $$
  \Big\langle\frac{D}{\partial t}\frac{D}{\partial s}V - \frac{D}{\partial s}\frac{D}{\partial t}V,\ W\Big\rangle
    = \Big\langle R\Big(\frac{\partial f}{\partial s},\frac{\partial f}{\partial t}\Big)V,\ W\Big\rangle
  $$
  for every $W\in T_{f(s_0,t_0)}M$.

  The covariant derivatives are \cref{def:dc-ch9-2-covariant-pair} pairs. The hypotheses
  constrain $\frac{D}{\partial s}V$, $\frac{D}{\partial t}V$ on the two-dimensional window
  (the outer covariant derivatives differentiate them along the transverse slice), and
  $\frac{D}{\partial t}\frac{D}{\partial s}V$, $\frac{D}{\partial s}\frac{D}{\partial t}V$ only
  on their own slice line.
\end{lemma}
\begin{proof}
  \leanok
  \uses{lem:dc-ch9-2-8-curvature-substitution, lem:dc-ch9-2-4-symmetry-manifold}
  \sketch
  Choose a common chart for the two slice curves. The covariant-pair hypotheses
  identify all four intrinsic derivatives with their coordinate counterparts, so
  \cref{lem:dc-ch9-2-8-curvature-substitution} applies. The coordinate equality
  transports to the intrinsic tangent space and yields the displayed formula.
\end{proof}
\begin{lemma}[proper second variation on a segment]
  \label{lem:dc-ch9-2-8-segment}
  \lean{Riemannian.Variation.deriv_deriv_dcEnergy_eq_second_variation, Riemannian.Variation.integral_second_variation_integrand_eq, Riemannian.Variation.curvatureFormAt_swap_pairs}
  \uses{lem:dc-ch9-2-8-energy-snd-deriv, lem:dc-ch9-2-4-symmetry-manifold, lem:dc-ch9-2-8-curvature-substitution-manifold, lem:dc-ch9-2-4-parts, lem:dc-ch9-2-5-geodesic, def:dc-ch9-2-2-energy, def:dc-ch9-2-covariant-pair}
  \leanok
  \dcref{ch9:2.8}
  Let $\gamma=f(s_0,\cdot):[a,b]\to M$ be a geodesic and $f$ a variation of $\gamma$ with
  variational field $V=\frac{\partial f}{\partial s}(s_0,\cdot)$ satisfying $V(a)=V(b)=0$ and
  $\frac{D}{\partial s}\frac{\partial f}{\partial s}(s_0,a)
    =\frac{D}{\partial s}\frac{\partial f}{\partial s}(s_0,b)=0$ (all of which hold when $f$ is
  proper). Then, on the segment $[a,b]$ carrying no breakpoint,
  $$
  \frac{1}{2}E''(s_0)=-\int_a^b\Big\langle V,\frac{D^2 V}{dt^2}+R(\gamma',V)\gamma'\Big\rangle dt .
  $$

\end{lemma}
\begin{proof}
  \leanok
  \uses{lem:dc-ch9-2-8-energy-snd-deriv, lem:dc-ch9-2-4-symmetry-manifold,
    lem:dc-ch9-2-8-curvature-substitution-manifold, lem:dc-ch9-2-4-parts, lem:dc-ch9-2-5-geodesic}
  \sketch
  By \cref{lem:dc-ch9-2-8-energy-snd-deriv},
  $\frac12 E''(s_0)=\int_a^b\{\langle\frac{D}{\partial s}\frac{D}{\partial s}\frac{\partial f}{\partial t},\frac{\partial f}{\partial t}\rangle+\langle\frac{D}{\partial s}\frac{\partial f}{\partial t},\frac{D}{\partial s}\frac{\partial f}{\partial t}\rangle\}\,dt$.
  The symmetry of the connection (\cref{lem:dc-ch9-2-4-symmetry-manifold}) rewrites the second
  integrand as $\langle V',V'\rangle$, which integration by parts
  (\cref{lem:dc-ch9-2-4-parts}, the variation being proper) turns into
  $-\int_a^b\langle V,\frac{D^2V}{dt^2}\rangle\,dt$. For the first integrand, the symmetry of
  the connection followed by the Ricci identity
  (\cref{lem:dc-ch9-2-8-curvature-substitution-manifold}) gives
  $\langle\frac{D}{\partial s}\frac{D}{\partial s}\frac{\partial f}{\partial t},\frac{\partial f}{\partial t}\rangle
    =\langle\frac{D}{\partial t}\frac{D}{\partial s}\frac{\partial f}{\partial s},\frac{\partial f}{\partial t}\rangle
     -R\big(\tfrac{\partial f}{\partial s},\tfrac{\partial f}{\partial t},\tfrac{\partial f}{\partial s},\tfrac{\partial f}{\partial t}\big)$,
  whose first term integrates to $0$ because $\gamma$ is a geodesic and
  $\frac{D}{\partial s}\frac{\partial f}{\partial s}$ vanishes at the endpoints
  (\cref{lem:dc-ch9-2-5-geodesic}), and whose curvature term is $-\langle R(\gamma',V)\gamma',V\rangle$:
  at $s=0$, $\frac{\partial f}{\partial s}=V$ and $\frac{\partial f}{\partial t}=\gamma'$, so it reads
  $R(V,\gamma',V,\gamma')$, which the interchange of the two antisymmetric pairs turns into
  $R(\gamma',V,\gamma',V)=\langle R(\gamma',V)\gamma',V\rangle$.
  Summing gives formula (3).
\end{proof}
\begin{proposition}[formula for the second variation]
  \label{prop:dc-ch9-2-8}
  \uses{def:dc-ch9-2-1, def:dc-ch9-2-2-energy, def:dc-ch9-2-covariant-pair, def:dc-ch4-2-1}
  \dcref{ch9:2.8}
  Let $\gamma:[0,a]\to M$ be a geodesic and let
  $f:(-\varepsilon,\varepsilon)\times[0,a]\to M$ be a proper variation of $\gamma$. Let $E$
  be the energy function of the variation. Then

  $$
  \frac{1}{2}E''(0)=-\int_0^a\Big\langle V(t),\frac{D^2 V}{dt^2}+R\Big(\frac{d\gamma}{dt},V\Big)\frac{d\gamma}{dt}\Big\rangle dt-\sum_{i=1}^k\Big\langle V(t_i),\frac{DV}{dt}(t_i^+)-\frac{DV}{dt}(t_i^-)\Big\rangle,\qquad(3)
  $$

  where $V$ is the variational field of $f$, $R$ is the curvature of $M$ and
  $\frac{DV}{dt}(t_i^+)=\lim_{t\to t_i,\,t>t_i}\frac{DV}{dt}$,
  $\frac{DV}{dt}(t_i^-)=\lim_{t\to t_i,\,t<t_i}\frac{DV}{dt}$.
\end{proposition}
\begin{proof}
  \uses{lem:dc-ch9-2-8-energy-snd-deriv, lem:dc-ch9-2-4-symmetry-manifold, lem:dc-ch9-2-8-curvature-substitution-manifold, lem:dc-ch9-2-4-parts, lem:dc-ch9-2-5-geodesic, lem:dc-ch9-2-8-jump-sum}
  \sketch
  Taking the derivative of (2),

  $$
  \frac{1}{2}\frac{d^2 E}{ds^2}=\sum_{i=0}^k\Big\langle\frac{D}{ds}\frac{\partial f}{\partial s},\frac{\partial f}{\partial t}\Big\rangle\Big|_{t_i}^{t_{i+1}}+\sum_{i=0}^k\Big\langle\frac{\partial f}{\partial s},\frac{D}{ds}\frac{\partial f}{\partial t}\Big\rangle\Big|_{t_i}^{t_{i+1}}-\int_0^a\Big\langle\frac{D}{ds}\frac{\partial f}{\partial s},\frac{D}{dt}\frac{\partial f}{\partial t}\Big\rangle dt-\int_0^a\Big\langle\frac{\partial f}{\partial s},\frac{D}{ds}\frac{D}{dt}\frac{\partial f}{\partial t}\Big\rangle dt.
  $$

  Putting $s=0$, the first and third terms are zero, since $f$ is proper and
  $\gamma$ is a geodesic. Because

  $$
  \frac{D}{ds}\frac{D}{dt}\frac{\partial f}{\partial t}=\frac{D}{dt}\frac{D}{ds}\frac{\partial f}{\partial t}+R\Big(\frac{\partial f}{\partial t},\frac{\partial f}{\partial s}\Big)\frac{\partial f}{\partial t},
  $$

  we have, at $s=0$,
  $\frac{D}{ds}\frac{D}{dt}\frac{\partial f}{\partial t}=\frac{D^2}{dt^2}V+R(\frac{d\gamma}{dt},V)\frac{d\gamma}{dt}$.
  Further use of the fact that the variation is proper yields

  $$
  \sum_{i=0}^k\Big\langle\frac{\partial f}{\partial s},\frac{D}{ds}\frac{\partial f}{\partial t}\Big\rangle\Big|_{t_i}^{t_{i+1}}=-\sum_{i=1}^k\Big\langle V(t_i),\frac{DV}{dt}(t_i^+)-\frac{DV}{dt}(t_i^-)\Big\rangle.\qquad(4)
  $$

  Putting all these facts together, we obtain (3).
\end{proof}
\begin{remark}[nonproper second variation with boundary terms]
  \label{rem:dc-ch9-2-9}
  \uses{prop:dc-ch9-2-8}
  \lean{Riemannian.Variation.hasDerivAt_deriv_dcEnergy_eq_piecewise_second_variation_nonproper,
    Riemannian.Variation.deriv_deriv_dcEnergy_eq_piecewise_second_variation_nonproper}
  \dcref{ch9:2.9}
  If the variation is not proper, the endpoint terms in the second-variation
  formula need not vanish, and the formula becomes

  $$
  \frac{1}{2}E''(0)=-\int_0^a\Big\langle V(t),\frac{D^2 V}{dt^2}+R\Big(\frac{d\gamma}{dt},V\Big)\frac{d\gamma}{dt}\Big\rangle dt-\sum_{i=1}^k\Big\langle V(t_i),\frac{DV}{dt}(t_i^+)-\frac{DV}{dt}(t_i^-)\Big\rangle
  $$

  $$
  -\Big\langle\frac{D}{ds}\frac{\partial f}{\partial s},\frac{d\gamma}{dt}\Big\rangle(0,0)+\Big\langle\frac{D}{ds}\frac{\partial f}{\partial s},\frac{d\gamma}{dt}\Big\rangle(0,a)-\Big\langle V(0),\frac{DV}{dt}(0)\Big\rangle+\Big\langle V(a),\frac{DV}{dt}(a)\Big\rangle.\qquad(5)
  $$
\end{remark}
\begin{lemma}[nonproper second variation on a segment]
  \label{lem:dc-ch9-2-9-segment}
  \lean{Riemannian.Variation.deriv_deriv_dcEnergy_eq_second_variation_nonproper, Riemannian.Variation.integral_second_variation_integrand_eq_nonproper}
  \uses{lem:dc-ch9-2-8-energy-snd-deriv, lem:dc-ch9-2-4-symmetry-manifold, lem:dc-ch9-2-8-curvature-substitution-manifold, lem:dc-ch9-2-4-parts, def:dc-ch9-2-2-energy, def:dc-ch9-2-covariant-pair}
  \leanok
  \dcref{ch9:2.9}
  Let $\gamma=f(s_0,\cdot):[a,b]\to M$ be a geodesic and $f$ a variation of $\gamma$ that need
  \emph{not} be proper. Then, on the segment $[a,b]$ carrying no breakpoint,
  $$
  \frac{1}{2}E''(s_0)=-\int_a^b\Big\langle V,\frac{D^2 V}{dt^2}+R(\gamma',V)\gamma'\Big\rangle dt
    +\Big\langle\frac{D}{\partial s}\frac{\partial f}{\partial s},\gamma'\Big\rangle\Big|_a^b
    +\big\langle V,V'\big\rangle\Big|_a^b .
  $$

\end{lemma}
\begin{proof}
  \leanok
  \uses{lem:dc-ch9-2-8-energy-snd-deriv, lem:dc-ch9-2-4-symmetry-manifold,
    lem:dc-ch9-2-8-curvature-substitution-manifold, lem:dc-ch9-2-4-parts}
  \sketch
  Identical to the assembly of \cref{lem:dc-ch9-2-8-segment}, except that the two integrations
  by parts (\cref{lem:dc-ch9-2-4-parts}) are applied in their boundary-keeping form: the
  $\langle V',V'\rangle$ term integrates to
  $\langle V,V'\rangle|_a^b-\int_a^b\langle V,\frac{D^2V}{dt^2}\rangle$, and the geodesic term
  $\langle\frac{D}{\partial t}\frac{D}{\partial s}\frac{\partial f}{\partial s},\gamma'\rangle$
  integrates to $\langle\frac{D}{\partial s}\frac{\partial f}{\partial s},\gamma'\rangle|_a^b$
  (the remaining integral vanishes as $\gamma$ is a geodesic). The curvature term is unchanged
  from \cref{lem:dc-ch9-2-8-segment}.
\end{proof}
\begin{remark}[the index form]
  \label{rem:dc-ch9-2-10}
  \lean{Riemannian.Variation.deriv_deriv_dcEnergy_eq_piecewise_indexForm, Riemannian.Variation.deriv_deriv_dcEnergy_eq_piecewise_indexForm_of_proper}
  \uses{prop:dc-ch9-2-8, rem:dc-ch9-2-9}
  \dcref{ch9:2.10}
  On each interval where $V$ is differentiable,
  $\frac{d}{dt}\langle V,\frac{DV}{dt}\rangle=\langle V,\frac{D^2 V}{dt^2}\rangle+\langle\frac{DV}{dt},\frac{DV}{dt}\rangle$.
  For a geodesic $\gamma:[0,a]\to M$ and a partition
  $0=t_0<t_1<\cdots<t_{k+1}=a$, formula (5) is

  $$
  \frac{1}{2}E''(0)=\int_0^a\{\langle V',V'\rangle-\langle R(\gamma',V)\gamma',V\rangle\}dt-\Big\langle\frac{D}{ds}\frac{\partial f}{\partial s},\gamma'\Big\rangle(0,0)+\Big\langle\frac{D}{ds}\frac{\partial f}{\partial s},\gamma'\Big\rangle(0,a).
  $$

  Define

  $$
  \int_0^a\{\langle V',V'\rangle-\langle R(\gamma',V)\gamma',V\rangle\}dt=I_a(V,V).
  $$

  If the variation is proper, $I_a(V,V)=\frac{1}{2}E''(0)$ depends only on $V$. In
  the general case,

  $$
  \frac{1}{2}E''(0)=I_a(V,V)-\Big\langle\frac{D}{ds}\frac{\partial f}{\partial s},\gamma'\Big\rangle(0,0)+\Big\langle\frac{D}{ds}\frac{\partial f}{\partial s},\gamma'\Big\rangle(0,a).\qquad(6)
  $$
\end{remark}
\begin{lemma}[differentiating $\langle V,DV/dt\rangle$]
  \label{lem:dc-ch9-2-10-differentiation}
  \lean{Riemannian.Variation.hasDerivAt_metricInner_covariantDeriv}
  \uses{lem:dc-ch9-2-covariant-pair-leibniz}
  \leanok
  \dcref{ch9:2.10}
  On each interval where $V$ is differentiable,
  $$
  \frac{d}{dt}\Big\langle V,\frac{DV}{dt}\Big\rangle
    =\Big\langle\frac{DV}{dt},\frac{DV}{dt}\Big\rangle
     +\Big\langle V,\frac{D^2V}{dt^2}\Big\rangle,
  $$
  the identity used to pass from (5) to (6).
\end{lemma}
\begin{proof}
  \leanok
  \uses{lem:dc-ch9-2-covariant-pair-leibniz}
  \sketch
  Apply the product rule \cref{lem:dc-ch9-2-covariant-pair-leibniz} to the pair $V$, with
  covariant derivative $\frac{DV}{dt}$, against the pair $\frac{DV}{dt}$, with covariant
  derivative $\frac{D^2V}{dt^2}$.
\end{proof}
\begin{definition}[the index form $I_a$]
  \label{def:dc-ch9-2-10-index-form}
  \lean{Riemannian.Variation.indexForm, Riemannian.Variation.indexForm_def}
  \uses{def:dc-ch9-2-covariant-pair, def:dc-ch4-2-1, def:dc-ch1-2-9}
  \leanok
  \dcref{ch9:2.10}
  For a geodesic $\gamma:[0,a]\to M$ and a vector field $V$ along $\gamma$ with covariant
  derivative $V'=\frac{DV}{dt}$, the \emph{index form} is
  $$
  I_a(V,V)=\int_0^a\{\langle V',V'\rangle-\langle R(\gamma',V)\gamma',V\rangle\}\,dt,
  $$
  with $R$ the curvature tensor and $\gamma'$ the geodesic velocity. For a
  proper variation, \cref{lem:dc-ch9-2-10-formula6} gives
  $I_a(V,V)=\frac{1}{2}E''(0)$.
\end{definition}
\begin{lemma}[second variation as the index form on a segment]
  \label{lem:dc-ch9-2-10-formula6}
  \lean{Riemannian.Variation.deriv_deriv_dcEnergy_eq_indexForm, Riemannian.Variation.integral_second_variation_eq_indexForm, Riemannian.Variation.ProperSecondVariationData.deriv_deriv_dcEnergy_eq_indexForm}
  \uses{def:dc-ch9-2-10-index-form, lem:dc-ch9-2-8-energy-snd-deriv, lem:dc-ch9-2-4-symmetry-manifold, lem:dc-ch9-2-8-curvature-substitution-manifold, lem:dc-ch9-2-5-geodesic}
  \leanok
  \dcref{ch9:2.10}
  Let $\gamma=f(s_0,\cdot):[a,b]\to M$ be a geodesic and $f$ a proper variation of $\gamma$
  (in the sense of \cref{lem:dc-ch9-2-8-segment}). Then
  $$
  \frac{1}{2}E''(0)=I_a(V,V).
  $$
\end{lemma}
\begin{proof}
  \leanok
  \uses{lem:dc-ch9-2-8-energy-snd-deriv, lem:dc-ch9-2-4-symmetry-manifold,
    lem:dc-ch9-2-8-curvature-substitution-manifold, lem:dc-ch9-2-5-geodesic,
    def:dc-ch9-2-10-index-form}
  \sketch
  By \cref{lem:dc-ch9-2-8-energy-snd-deriv},
  $\frac12 E''(0)=\int_a^b\{\langle\frac{D}{\partial s}\frac{D}{\partial s}\frac{\partial f}{\partial t},\frac{\partial f}{\partial t}\rangle+\langle\frac{D}{\partial s}\frac{\partial f}{\partial t},\frac{D}{\partial s}\frac{\partial f}{\partial t}\rangle\}\,dt$.
  The symmetry of the connection (\cref{lem:dc-ch9-2-4-symmetry-manifold}) makes the second
  integrand $\langle V',V'\rangle$; the symmetry followed by the Ricci identity
  (\cref{lem:dc-ch9-2-8-curvature-substitution-manifold}), together with the geodesic
  vanishing of the boundary term (\cref{lem:dc-ch9-2-5-geodesic}), makes the first integrand
  integrate to $-\int_a^b\langle R(\gamma',V)\gamma',V\rangle\,dt$. Their sum is $I_a(V,V)$.

\end{proof}
\begin{remark}[Index form as the Hessian of energy]
  \label{rem:dc-ch9-2-12}
  \uses{rem:dc-ch9-2-7}
  \dcref{ch9:2.12}
  At a geodesic $\gamma\in\Omega_{p,q}$, the quadratic form
  $2I_a(V,V)$ is the Hessian $d^2E(V,V)$ of the energy functional in the
  direction of the endpoint-vanishing field $V$.
\end{remark}
\section{The theorems of Bonnet--Myers and of Synge--Weinstein}

\begin{lemma}[a velocity-seeded parallel orthonormal frame along a geodesic]
  \label{lem:dc-ch9-3-1-velocity-frame}
  \lean{Riemannian.Variation.exists_velocitySeededParallelOrthoFrameAlongOn}
  \leanok
  \dcref{ch9:3.1}
  Let $\gamma:[a,b]\to M$ be a geodesic with $\gamma'(a)\ne 0$. Then there are parallel fields
  $e_1,\dots,e_n$ along $\gamma$, orthonormal at every $t\in[a,b]$, and a distinguished index
  $n_0$ with $e_{n_0}(t)=\ell^{-1}\gamma'(t)$ for all $t$, where $\ell=|\gamma'|>0$ is the
  (constant) speed. In particular $e_i(t)\perp\gamma'(t)$ for $i\ne n_0$.
\end{lemma}
\begin{proof}
  \leanok
  \sketch
  Extend the unit velocity $\ell^{-1}\gamma'(a)$ to a $g$-orthonormal basis of $T_{\gamma(a)}M$
  and parallel-transport it along $\gamma$; parallel transport is a linear isometry, so the frame
  stays orthonormal. The distinguished transported member and $\ell^{-1}\gamma'$ are both parallel
  fields with the same value at $a$, so by uniqueness of parallel transport they agree on
  $[a,b]$.
\end{proof}
\begin{lemma}[the index form of a scaled parallel field]
  \label{lem:dc-ch9-3-1-index-form-scaled}
  \lean{Riemannian.Variation.indexForm_smul_eq}
  \uses{def:dc-ch9-2-10-index-form, def:dc-ch9-2-covariant-pair}
  \leanok
  \dcref{ch9:3.1}
  For a field $V(t)=\varphi(t)\,e(t)$ with covariant derivative $V'=\varphi'\,e$ (as when $e$ is
  parallel), the index form is
  $$I_a(V,V)=\int_a^b\Big((\varphi')^2\,\langle e,e\rangle
    -\varphi^2\,\langle R(\gamma',e)\gamma',e\rangle\Big)\,dt.$$
  In particular, this applies to $V_j=(\sin\pi t)e_j$ in
  \cref{thm:dc-ch9-3-1}.
\end{lemma}
\begin{proof}
  \leanok
  \sketch
  Pull $\varphi'$ out of $\langle\varphi' e,\varphi' e\rangle$ (bilinearity of the metric) and
  $\varphi$ out of each $e$-argument of $\langle R(\gamma',\varphi e)\gamma',\varphi e\rangle$
  (homogeneity of the curvature form in slots 2 and 4).
\end{proof}
\begin{lemma}[the index form of $V_j$ in the velocity-seeded frame]
  \label{lem:dc-ch9-3-1-index-form-frame}
  \lean{Riemannian.Variation.indexForm_smul_frame_eq}
  \uses{lem:dc-ch9-3-1-index-form-scaled, lem:dc-ch9-3-1-velocity-frame}
  \leanok
  \dcref{ch9:3.1}
  With the frame of \cref{lem:dc-ch9-3-1-velocity-frame} (so $e=e_j$ is a unit field and
  $\gamma'=\ell\,e_n$), for $V=\varphi\,e_j$
  $$I_a(V,V)=\int_a^b\Big((\varphi')^2-\varphi^2\,\ell^2\,\langle R(e_n,e_j)e_n,e_j\rangle\Big)\,dt.$$
  For $\varphi=\sin\pi t$ and orthonormal $e_n\perp e_j$, the curvature
  term is $\ell^2\sin^2(\pi t)K(e_n,e_j)$.
\end{lemma}
\begin{proof}
  \leanok
  \uses{lem:dc-ch9-3-1-index-form-scaled}
  \sketch
  Substitute $\langle e_j,e_j\rangle=1$ and $\gamma'=\ell\,e_n$ into
  \cref{lem:dc-ch9-3-1-index-form-scaled} and use homogeneity of the curvature form in slots
  1 and 3.
\end{proof}
\begin{lemma}[the proper sine exponential variation]
  \label{lem:dc-ch9-3-1-sine-variation}
  \lean{Riemannian.Variation.bonnetMyersSineVariation, Riemannian.Variation.bonnetMyersSineVariation_zero, Riemannian.Variation.bonnetMyersSineVariation_left, Riemannian.Variation.bonnetMyersSineVariation_right, Riemannian.Variation.hasDerivAt_bonnetMyersSineVariation_zero, Riemannian.Variation.exists_contMDiffOn_infty_bonnetMyersSineVariation_strip, Riemannian.Variation.exists_proper_sine_smoothExponentialVariation, Riemannian.Variation.exists_proper_sine_exponentialVariation}
  \leanok
  \dcref{ch9:3.1}
  Let $\gamma:\mathbb R\to M$ be a geodesic in a complete Riemannian
  manifold, and let $e$ be a parallel field along an open interval containing
  $[0,1]$. Then
  \[
    f(s,t)=\exp_{\gamma(t)}\!\bigl(s\sin(\pi t)e(t)\bigr)
  \]
  is smooth on an open product neighborhood of
  $\{0\}\times[0,1]$, satisfies $f(0,t)=\gamma(t)$, fixes both endpoints, and
  has variational field $V(t)=\sin(\pi t)e(t)$.
\end{lemma}
\begin{proof}
  \leanok
  \sketch
  Completeness makes every exponential map defined on the whole tangent
  space. Smooth dependence of geodesics on their initial data, together with
  smoothness of the parallel field, gives a common smooth product
  neighborhood of the compact zero slice. At $s=0$ the exponential of the
  zero vector is $\gamma(t)$; at $t=0,1$ the sine factor vanishes. Finally,
  differentiating $s\mapsto\exp_{\gamma(t)}(sV(t))$ at zero gives $V(t)$.
\end{proof}
\begin{lemma}[energy minimality of the sine variation]
  \label{lem:dc-ch9-3-1-sine-energy-min}
  \lean{Riemannian.Variation.isLocalMin_dcEnergy_bonnetMyersSineVariation}
  \uses{lem:dc-ch9-2-3-smooth-variation-energy-min, lem:dc-ch9-3-1-sine-variation}
  \leanok
  \dcref{ch9:3.1}
  If $\gamma:[0,1]\to M$ is a distance-minimizing geodesic and $e$ is a
  parallel field as in \cref{lem:dc-ch9-3-1-sine-variation}, then the energy
  of the sine exponential variation has a local minimum at $s=0$.
\end{lemma}
\begin{proof}
  \leanok
  \sketch
  The sine exponential variation is smooth on a product neighborhood, has
  zero slice $\gamma$, and fixes both endpoints by
  \cref{lem:dc-ch9-3-1-sine-variation}. Apply
  \cref{lem:dc-ch9-2-3-smooth-variation-energy-min}.
\end{proof}
\begin{lemma}[second-order necessary condition at a minimum]
  \label{lem:dc-ch9-3-1-second-order-min}
  \lean{Riemannian.Variation.isLocalMin_deriv_deriv_nonneg, Riemannian.Variation.indexForm_nonneg_of_energyLocalMin, Riemannian.Variation.indexForm_nonneg_of_energyLocalMin_at}
  \leanok
  \dcref{ch9:3.1}
  If $f:\mathbf{R}\to\mathbf{R}$ has a local minimum at $x$, is continuous there and is twice
  differentiable at $x$, then $f''(x)\ge 0$. This is the step by which a \emph{minimizing}
  geodesic forces $E_j''(0)\ge 0$ in \cref{thm:dc-ch9-3-1}: each curve of the variation joins
  the same endpoints, so $E_j(s)\ge E_j(0)$, i.e. $s=0$ is a minimum of $E_j$.
\end{lemma}
\begin{proof}
  \leanok
  \sketch
  If $f''(x)<0$, then for sufficiently small $h\ne0$ the symmetric second
  difference $f(x+h)+f(x-h)-2f(x)$ is negative. One of $f(x+h)$ and
  $f(x-h)$ is therefore smaller than $f(x)$, contradicting local minimality.
\end{proof}
\begin{lemma}[the summed index form is negative]
  \label{lem:dc-ch9-3-1-index-sum-neg}
  \lean{Riemannian.Variation.integral_frame_sum_lt_zero, Riemannian.Variation.sum_indexForm_smul_frame_neg, Riemannian.Variation.exists_indexForm_smul_frame_neg}
  \uses{lem:dc-ch9-3-1-velocity-frame, lem:dc-ch9-3-1-index-form-frame, lem:dc-ch9-3-3-ricci-sectional, def:dc-ch9-2-10-index-form}
  \leanok
  \dcref{ch9:3.1}
  Let $\gamma:[0,1]\to M$ be a geodesic of constant speed $\ell$ with the velocity-seeded
  parallel orthonormal frame of \cref{lem:dc-ch9-3-1-velocity-frame} (so each $e_j$, $j\ne n_0$,
  is a unit field and $\gamma'=\ell\,e_{n_0}$), and put $V_j(t)=(\sin\pi t)\,e_j(t)$. If $n\ge 2$
  (there is at least one index $j\ne n_0$), $r>0$, $\pi r<\ell$, and the raw curvature sum satisfies
  $\sum_{j\ne n_0}\langle R(e_{n_0},e_j)e_{n_0},e_j\rangle\ge (n-1)/r^2$ at every $t\in[0,1]$
  (each term continuous in $t$), then
  $$
  \sum_{j\ne n_0} I_a(V_j,V_j)<0,\qquad\text{hence}\quad I_a(V_j,V_j)<0\ \text{for some }j\ne n_0.
  $$

  By \cref{lem:dc-ch9-3-3-ricci-sectional}, the displayed curvature sum is
  the unnormalized Ricci form $Q(e_{n_0},e_{n_0})$.
\end{lemma}
\begin{proof}
  \leanok
  \uses{lem:dc-ch9-3-1-index-form-frame}
  \sketch
  For each $j\ne n_0$, \cref{lem:dc-ch9-3-1-index-form-frame} with $\varphi=\sin\pi t$
  (so $\varphi'=\pi\cos\pi t$) gives
  $I_a(V_j,V_j)=\int_0^1\big((\pi\cos\pi t)^2-\sin^2\pi t\,(\ell^2\langle R(e_{n_0},e_j)e_{n_0},e_j\rangle)\big)\,dt$.
  The integrands are continuous, so the finite sum commutes with the integral, and summing over
  the $n-1$ indices $j\ne n_0$ collapses it to
  $$
  \int_0^1\Big((n-1)(\pi\cos\pi t)^2-\sin^2\pi t\,\big(\ell^2\textstyle\sum_{j\ne n_0}\langle R(e_{n_0},e_j)e_{n_0},e_j\rangle\big)\Big)\,dt.
  $$
  Write $S(t)=\ell^2\sum_{j\ne n_0}\langle R(e_{n_0},e_j)e_{n_0},e_j\rangle$. The integrand equals
  $(n-1)\pi^2\cos 2\pi t-\sin^2\pi t\,(S(t)-(n-1)\pi^2)$; the first part integrates to $0$
  (as $\int_0^1\cos 2\pi t\,dt=0$), and the second is the integral of a function that is $\ge 0$
  on $[0,1]$ and strictly positive on $(0,1)$ --- there $\sin\pi t>0$ and, since $\pi r<\ell$
  gives $\pi^2 r^2<\ell^2$, $S(t)\ge\ell^2(n-1)/r^2>(n-1)\pi^2$. Hence the whole integral is
  $<0$. A strictly negative finite sum has a strictly negative summand.
\end{proof}
\begin{theorem}[Bonnet--Myers]
  \label{thm:dc-ch9-3-1}
  \uses{lem:dc-ch9-2-3, prop:dc-ch9-2-8, lem:dc-ch9-3-1-velocity-frame, lem:dc-ch9-3-1-index-form-frame, lem:dc-ch9-3-1-index-sum-neg, lem:dc-ch9-3-1-sine-variation, lem:dc-ch9-3-1-sine-energy-min, lem:dc-ch9-3-1-second-order-min, lem:dc-ch9-3-3-ricci-sectional}
  \notready
  \dcref{ch9:3.1}
  \notready
  Let $M^n$ be a complete Riemannian manifold. Suppose that the Ricci curvature of
  $M$ satisfies

  $$
  \mathrm{Ric}_p(v)\ge\frac{1}{r^2}>0,
  $$

  for all $p\in M$ and for all $v\in T_p(M)$. Then $M$ is compact and the diameter
  $\mathrm{diam}(M)\le\pi r$.
\end{theorem}
\begin{proof}
  \sketch
  Let $p$ and $q$ be any two points in $M$. Since $M$ is complete, there
  exists a minimizing geodesic $\gamma:[0,1]\to M$ joining $p$ to $q$. It is enough
  to show that the length $\ell(\gamma)\le\pi r$, because then $M$ is bounded and
  complete, therefore compact; and since $d(p,q)\le\pi r$ for any $p,q$, it follows
  that $\mathrm{diam}(M)\le\pi r$. Suppose, to the contrary, that
  $\ell(\gamma)=\ell>\pi r$. Consider parallel fields $e_1(t),\dots,e_{n-1}(t)$ along
  $\gamma$ which are orthonormal for each $t\in[0,1]$ and belong to the orthogonal
  complement of $\gamma'(t)$. Let $e_n(t)=\frac{\gamma'(t)}{\ell}$ and let $V_j$ be
  the vector field along $\gamma$ given by $V_j(t)=(\sin\pi t)e_j(t)$,
  $j=1,\dots,n-1$. Then $V_j(0)=V_j(1)=0$, so $V_j$ generates a proper variation of
  $\gamma$ with energy $E_j$. Using the formula for the second variation and the fact
  that $e_j$ is parallel,

  $$
  \frac{1}{2}E_j''(0)=-\int_0^1\langle V_j,V_j''+R(\gamma',V_j)\gamma'\rangle dt=\int_0^1\sin^2\pi t\,(\pi^2-\ell^2 K(e_n(t),e_j(t)))\,dt,
  $$

  where $K(e_n(t),e_j(t))$ is the sectional curvature at $\gamma(t)$ with respect to
  the plane generated by $e_n(t),e_j(t)$. Summing on $j$ and using the definition of
  Ricci curvature,

  $$
  \frac{1}{2}\sum_{j=1}^{n-1}E_j''(0)=\int_0^1\{\sin^2\pi t\,((n-1)\pi^2-(n-1)\ell^2\,\mathrm{Ric}_{\gamma(t)}(e_n(t)))\}\,dt.
  $$

  Since $\mathrm{Ric}_{\gamma(t)}(e_n(t))\ge\frac{1}{r^2}$ and $\ell>\pi r$, we have
  $(n-1)\ell^2\,\mathrm{Ric}_{\gamma(t)}(e_n(t))>(n-1)\pi^2$, hence

  $$
  \frac{1}{2}\sum_{j=1}^{n-1}E_j''(0)<\int_0^1\sin^2\pi t\,((n-1)\pi^2-(n-1)\pi^2)\,dt=0.
  $$

  Thus some $E_j''(0)<0$, which, by \cref{lem:dc-ch9-2-3},
  contradicts the fact that $\gamma$ is minimizing. Therefore $\ell\le\pi r$.
\end{proof}
\begin{corollary}[finiteness of the fundamental group]
  \label{cor:dc-ch9-3-2}
  \dcref{ch9:3.2}
  Let $M$ be a complete Riemannian manifold with $\mathrm{Ric}_p(v)\ge\delta>0$, for
  all $p\in M$ and all $v\in T_p(M)$. Then the universal cover of $M$ is compact. In
  particular, the fundamental group $\pi_1(M)$ is finite.
\end{corollary}

\begin{lemma}[finite fundamental group from a compact simply connected cover]
  \label{lem:dc-ch9-3-2-compact-cover}
  \dcref{ch9:3.2}
  \lean{Riemannian.Variation.finite_fiber_of_compact_covering,
    Riemannian.Variation.monodromy_at_injective,
    Riemannian.Variation.finite_fundamentalGroup_of_compact_simplyConnected_cover}
  \leanok
  Let $p:X\to Y$ be a surjective covering map, with $X$ compact and simply connected and
  $Y$ a $T_1$ space.  Then the fundamental group $\pi_1(Y,y)$ is finite for every basepoint
  $y\in Y$.
\end{lemma}
\begin{proof}
  \leanok
  A covering trivialization makes each fibre discrete, while compactness makes it finite.
  Simple connectedness of $X$ makes monodromy at a chosen point of the fibre injective on
  $\pi_1(Y,y)$; hence the fundamental group injects into a finite fibre.
\end{proof}

\begin{proof}
  \leanok
  \sketch
  Introducing on the universal cover $\pi:\tilde M\to M$ the covering metric
  (the metric such that $\pi$ is a local isometry), $\tilde M$ is complete and its
  Ricci curvature satisfies $\widetilde{\mathrm{Ric}}_p\ge\delta>0$. From the theorem,
  $\tilde M$ is compact. Hence the number of sheets of the covering is finite; since
  that is the number of elements in $\pi_1(M)$, we conclude that $\pi_1(M)$ is
  finite.
\end{proof}
\begin{lemma}[Ricci as a sum of sectional curvatures]
  \label{lem:dc-ch9-3-3-ricci-sectional}
  \lean{Riemannian.sectionalCurvature_eq_of_orthonormal, Riemannian.ricciForm_self_eq_sum_sectionalCurvature, Riemannian.ricciForm_self_ge_of_sectionalCurvature_ge, Riemannian.Variation.HasSectionalCurvatureLowerBound, Riemannian.Variation.hasRicciLowerBound_of_sectionalCurvatureLowerBound}
  \uses{def:dc-ch4-3-2, def:dc-ch4-4-ricci-form}
  \leanok
  \dcref{ch9:3.3}
  Let $\{e_1,\dots,e_n\}$ be an orthonormal basis of $T_pM$ and put $v=e_{i_0}$. Then
  $$
  Q(v,v)=\sum_{i\ne i_0}K(v,e_i),
  $$
  where $Q$ is the (unnormalized) Ricci form and $K$ the sectional curvature.
  Consequently, if $K(v,e_i)\ge k$ for every $i\ne i_0$, then $Q(v,v)\ge(n-1)k$.

  With the normalized Ricci curvature
  $\mathrm{Ric}_p(v)=Q(v,v)/(n-1)$, the same hypothesis gives
  $\mathrm{Ric}_p(v)\ge k$. The sum runs over $n-1$ terms because the diagonal contribution $B(v,v,v,v)$
  vanishes, and each pair $(v,e_i)$, $i\ne i_0$, being orthonormal, has $|v\wedge e_i|^2=1$,
  so its sectional curvature is the bare numerator $B(v,e_i,v,e_i)$.
\end{lemma}
\begin{proof}
  \leanok
  \uses{def:dc-ch4-3-2, def:dc-ch4-4-ricci-form}
  \sketch
  By the orthonormal-basis formula for the Ricci form, $Q(v,v)=\sum_i B(v,e_i,v,e_i)$. The
  term $i=i_0$ is $B(v,v,v,v)=0$ by the antisymmetry of $B$ in its first two slots. For
  $i\ne i_0$ the pair $(v,e_i)$ is orthonormal, so
  $|v\wedge e_i|^2=\langle v,v\rangle\langle e_i,e_i\rangle-\langle v,e_i\rangle^2=1$ and
  hence $K(v,e_i)=B(v,e_i,v,e_i)/1=B(v,e_i,v,e_i)$. Summing gives the identity. The
  inequality follows by bounding each of the $n-1$ summands below by $k$.
\end{proof}
\begin{corollary}[positive sectional curvature bound]
  \label{cor:dc-ch9-3-3}
  \uses{thm:dc-ch9-3-1, cor:dc-ch9-3-2, lem:dc-ch9-3-3-ricci-sectional}
  \dcref{ch9:3.3}
  Let $M$ be a complete Riemannian manifold with sectional curvature
  $K\ge\frac{1}{r^2}>0$. Then $M$ is compact, $\mathrm{diam}(M)\le\pi r$ and
  $\pi_1(M)$ is finite.
\end{corollary}
\begin{remark}[Positivity without a uniform lower bound]
  \label{rem:dc-ch9-3-4}
  \dcref{ch9:3.4}
  The hypothesis $K\ge\delta>0$ cannot be relaxed to $K>0$. Indeed, the paraboloid
  $\{(x,y,z)\in\mathbf{R}^3;z=x^2+y^2\}$ has curvature $K>0$, but is complete and
  non-compact.
\end{remark}
\begin{remark}[Sharpness of the Bonnet--Myers diameter bound]
  \label{rem:dc-ch9-3-6}
  \uses{thm:dc-ch9-3-1}
  \dcref{ch9:3.6}
  The estimate for the diameter given by \cref{thm:dc-ch9-3-1} cannot be improved. Indeed, the
  unit sphere $S^n\subset\mathbf{R}^{n+1}$ has constant sectional curvature equal to
  1 (therefore its Ricci curvature also is a constant equal to 1) and
  $\mathrm{diam}(S^n)=\pi$. This example is unique in the following sense: let $M^n$
  be complete with $\mathrm{Ric}_p(v)\ge 1/r^2$, for all $p\in M$ and all $v\in T_pM$;
  if $\mathrm{diam}(M)=\pi r$, then $M^n$ is isometric to the sphere $S^n$ of
  curvature $1/r^2$.
\end{remark}
\begin{theorem}[Synge--Weinstein]
  \label{thm:dc-ch9-3-7}
  \dcref{ch9:3.7}
  Let $f$ be an isometry of a compact oriented Riemannian manifold $M^n$. Suppose
  that $M$ has positive sectional curvature and that $f$ preserves the orientation of
  $M$ if $n$ is even, and reverses it if $n$ is odd. Then $f$ has a fixed point, i.e.,
  there exists $p\in M$ with $f(p)=p$.
\end{theorem}
\begin{proof}
  \proves{thm:dc-ch9-3-7}
  \uses{rem:dc-ch9-2-9, lem:dc-ch9-3-8, lem:dc-ch9-2-2-schwarz}
  \sketch
  Suppose, to the contrary, that $f(q)\ne q$ for all $q\in M$.
  Let $p\in M$ such that $d(p,f(p))$ attains a minimum. Since $M$ is complete, there
  exists a normalized minimizing geodesic $\gamma:[0,\ell]\to M$, joining $p$ to
  $f(p)$, that is, $\gamma(0)=p$, $\gamma(\ell)=f(p)$. Let
  $\tilde A=P\circ df_p:T_p(M)\to T_p(M)$, where $P$ is the parallel transport along
  $\gamma$ from $f(p)=\gamma(\ell)$ to $p$. Then $\tilde A$ is an isometry. We show
  that $(f\circ\gamma)'(0)=\gamma'(\ell)$. Consider the geodesic $f\circ\gamma$ which
  joins $f(p)$ to $f^2(p)$. Let $p'=\gamma(t')$, $t'\ne 0$, $t'\ne\ell$, and
  $f(p')=f\circ\gamma(t')$. Since $f$ is an isometry, $d(p,p')=d(f(p),f(p'))$ and,
  from the triangle inequality,

  $$
  d(p',f(p'))\le d(p',f(p))+d(f(p),f(p'))=d(p',f(p))+d(p,p')=d(p,f(p)).
  $$

  Because $d(p,f(p))$ is a minimum, $d(p',f(p'))=d(p',f(p))+d(f(p),f(p'))$. Therefore
  the curve formed by $\gamma$ and $f\circ\gamma$ is a geodesic, hence
  $(f\circ\gamma)'(0)=\gamma'(\ell)$, as asserted. It follows that $\tilde A$ leaves
  $\gamma'(0)$ fixed, since

  $$
  \tilde A(\gamma'(0))=P\circ df_p(\gamma'(0))=P((f\circ\gamma)'(0))=P(\gamma'(\ell))=\gamma'(0).
  $$

  Let $A$ be the restriction of $\tilde A$ to the orthogonal complement of
  $\gamma'(0)$. Then $A$ is orthogonal on $\mathbf{R}^{n-1}$ and, since $P$ is an
  orientation-preserving isometry,

  $$
  \det A=\det\tilde A=\det(P\circ df_p)=(-1)^n,
  $$

  where in the last equality we use the hypothesis on $f$ and the fact that $P$
  preserves orientation. By \cref{lem:dc-ch9-3-8}, $A$ leaves a vector invariant. Let $e_1(t)$
  be a unit parallel field along $\gamma$ such that, for each $t$, $e_1(t)$ belongs to
  the orthogonal complement of $\gamma'(t)$ and $e_1(0)$ is invariant by $A$. Let
  $\beta(s)$, $s\in(-\varepsilon,\varepsilon)$, be a geodesic such that $\beta(0)=p$ and
  $\beta'(0)=e_1(0)$. Because $P\circ df_p(e_1(0))=e_1(0)$, we have
  $df_p(e_1(0))=e_1(\ell)$, that is, the geodesic $f\circ\beta$ is such that
  $f\circ\beta(0)=f(p)$ and $(f\circ\beta)'(0)=e_1(\ell)$. Let $h$ be a variation of
  $\gamma$ given by $h(s,t)=\exp_{\gamma(t)}(se_1(t))$,
  $s\in(-\varepsilon,\varepsilon)$, $t\in[0,\ell]$. Since $h(s,0)=\beta(s)$, then
  $h(s,\ell)=\exp_{f(p)}(se_1(\ell))=(f\circ\beta)(s)$. Therefore,

  $$
  V(t)=\frac{\partial}{\partial s}\exp_{\gamma(t)}(se_1(t))\Big|_{s=0}=e_1(t),
  $$

  hence $\frac{D^2 V}{dt^2}=0$. The endpoint transverse curves are the
  geodesics $\beta$ and $f\circ\beta$, so their acceleration boundary terms vanish.
  Using the second-variation formula \cref{rem:dc-ch9-2-9} and
  $\frac{\partial h}{\partial s}(0,t)=e_1(t)$,

  $$
  \frac{1}{2}\frac{d^2 E}{ds^2}(0)=-\int_0^\ell\Big\langle V(t),R\Big(\frac{d\gamma}{dt},V\Big)\frac{d\gamma}{dt}\Big\rangle dt+\Big\langle\frac{d\gamma}{dt},\frac{D}{ds}\frac{\partial h}{\partial s}\Big\rangle(0,\ell)-\Big\langle\frac{d\gamma}{dt},\frac{D}{ds}\frac{\partial h}{\partial s}\Big\rangle(0,0)=-\int_0^\ell K\Big(e_1(t),\frac{d\gamma}{dt}\Big)dt.
  $$

  Because the sectional curvature is positive, $\frac{1}{2}\frac{d^2 E}{ds^2}(0)<0$,
  and therefore $\frac{dE}{ds}$ is strictly decreasing on a neighborhood of zero.
  This shows that there exists a curve $c$ in the variation such that
  $(\ell(c))^2\le\ell E(c)<\ell E(\gamma)=\ell(\gamma)^2$. Since the curves in the
  variation join $q$ to $f(q)$, we obtain a contradiction to the fact that
  $d(p,f(p))$ is a minimum.
\end{proof}
\begin{lemma}[invariant vector of an orthogonal map]
  \label{lem:dc-ch9-3-8}
  \lean{Riemannian.Variation.exists_fixed_vector_of_det_eq, Riemannian.Variation.exists_fixed_vector_of_orientation_parity}
  \leanok
  \dcref{ch9:3.8}
  Let $A$ be an orthogonal linear transformation of $\mathbf{R}^{n-1}$ and suppose
  that $\det A=(-1)^n$. Then $A$ fixes a non-zero vector of
  $\mathbf{R}^{n-1}$.
\end{lemma}
\begin{proof}
  \uses{rem:dc-ch9-2-9, lem:dc-ch9-3-7-orientation-det, lem:dc-ch9-3-8, lem:dc-ch9-2-2-schwarz}
  \sketch
  Since $A$ is orthogonal, $A-I=A(I-A^{T})$. Taking determinants in dimension
  $n-1$ gives
  $$
  \det(A-I)=\det(A)(-1)^{n-1}\det(A-I)=-\det(A-I).
  $$
  Hence $\det(A-I)=0$, so $1$ is an eigenvalue of $A$.
\end{proof}
\begin{remark}[Conformal Synge--Weinstein theorem]
  \label{rem:dc-ch9-3-9}
  \dcref{ch9:3.9}
  The conclusion of \cref{thm:dc-ch9-3-7} remains valid when $f$ is a
  conformal diffeomorphism. Extending it to arbitrary diffeomorphisms would
  preclude positive sectional curvature on $S^2\times S^2$, because the
  product of the antipodal maps is orientation-preserving and fixed-point free.
\end{remark}
\begin{corollary}[Synge]
  \label{cor:dc-ch9-3-10}
  \uses{thm:dc-ch9-3-7, cor:dc-ch9-3-3}
  \dcref{ch9:3.10}
  Let $M^n$ be a compact manifold with positive sectional curvature.
  \emph{(a)} If $M^n$ is orientable and $n$ is even, then $M$ is simply connected.
  \emph{(b)} If $n$ is odd, then $M^n$ is orientable.
\end{corollary}
\begin{proof}
  \sketch
  \emph{(a)} Equip the universal cover $\pi:\tilde M\to M$ with the
  covering metric and orientation. Compactness and positivity give a uniform
  lower curvature bound; hence \cref{cor:dc-ch9-3-3} makes $\tilde M$
  compact. Every deck transformation is an orientation-preserving isometry, so
  \cref{thm:dc-ch9-3-7} gives it a fixed point. A deck transformation with a
  fixed point is the identity, and therefore $\pi_1(M)$ is trivial.

  \emph{(b)} If $M$ were non-orientable, its orientable double cover
  $\overline M$ would be compact and carry the covering metric. Its nontrivial
  deck transformation is an orientation-reversing isometry without fixed
  points, contradicting \cref{thm:dc-ch9-3-7} when $n$ is odd.
\end{proof}
\begin{remark}[Projective-space counterexamples]
  \label{rem:dc-ch9-3-11}
  \dcref{ch9:3.11}
  The real projective space $P^2(\mathbf{R})$ of dimension two, which is compact,
  non-orientable and not simply connected, is an example which shows the necessity of
  orientability in a) and odd dimension in b). On the other hand, $P^3(\mathbf{R})$
  which is compact, orientable and not simply connected is an example which shows the
  necessity of even dimension in (a).
\end{remark}

\begin{lemma}[complete-exponential realization of a prescribed field]
  \label{lem:dc-ch9-2-2-exponential}
  \dcref{ch9:2.2}
  \lean{Riemannian.Variation.exponentialVariation,
    Riemannian.Variation.exponentialVariation_zero,
    Riemannian.Variation.exponentialVariation_isProperVariation,
    Riemannian.Variation.hasDerivAt_extChartAt_exponentialVariation_zero,
    Riemannian.Variation.variationalField_exponentialVariation,
    Riemannian.Variation.exponentialVariation_isVariationOfOrder_of_contMDiffOn,
    Riemannian.Variation.exponentialVariation_isVariation_of_contMDiffOn,
    Riemannian.Variation.exponentialVariation_isVariationOfOrder_of_piecewise_contMDiffOn,
    Riemannian.Variation.exponentialVariation_isVariation_of_piecewise_contMDiffOn,
    Riemannian.Variation.exists_variation_with_variationalField_of_contMDiffOn,
    Riemannian.Variation.exists_properVariation_with_variationalField_of_contMDiffOn,
    Riemannian.Variation.exists_properVariation_with_variationalField_of_piecewise_contMDiffOn}
  \leanok
  \uses{def:dc-ch9-2-1}
  Let $c:\mathbb R\to M$ and $V:\mathbb R\to TM$ be represented in the model
  vector space of a complete Riemannian manifold.  If, for some $a,\varepsilon>0$,
  the complete-exponential family
  $$f(s,t)=\exp_{c(t)}(sV(t))$$
  is $C^1$ on $(-\varepsilon,\varepsilon)\times[0,a]$, then $f$ is a variation of
  $c$ and its variational field is $V$.  If $V(0)=V(a)=0$, the same family is a
  proper variation.
\end{lemma}
\begin{proof}
  \leanok
  The zero-slice identity follows from $\exp_p(0)=p$.  Differentiating the radial
  exponential at the zero vector gives
  $\partial_s f(0,t)=V(t)$, so the variational field is prescribed.  The assumed
  $C^1$ regularity supplies the variation condition on the one-piece subdivision;
  when the endpoint values of $V$ vanish, the endpoint exponential identities make
  both endpoint curves constant, which is precisely properness.
\end{proof}

\begin{lemma}[equality gives pointwise proportional arc length]
  \label{lem:dc-ch9-2-3-pointwise-arclength}
  \dcref{ch9:2.3}
  \lean{Riemannian.pathELength_eq_of_dcEnergy_eq}
  \leanok
  \uses{lem:dc-ch9-2-3-equality-case, lem:dc-ch9-2-3-arclength-bridge}
  In the situation of \cref{lem:dc-ch9-2-3-equality-case}, for every $t\in[a,b]$,
  $$
  L\big(c|_{[a,t]}\big)=\frac{L(c)}{b-a}(t-a).
  $$
  Thus the parameter of $c$ is proportional to arc length pointwise, not merely almost
  everywhere.
\end{lemma}
\begin{proof}
  \leanok
  \uses{lem:dc-ch9-2-3-equality-case, lem:dc-ch9-2-3-arclength-bridge}
  By \cref{lem:dc-ch9-2-3-equality-case}, the speed of $c$ equals
  $L(c)/(b-a)$ almost everywhere on $(a,b]$. Restricting this equality to $(a,t]$ and
  integrating gives
  $L(c|_{[a,t]})=(L(c)/(b-a))(t-a)$. The identification of this integral with path length
  is \cref{lem:dc-ch9-2-3-arclength-bridge}. $\square$
\end{proof}

\begin{lemma}[finite-segment assembly of the first variation]
  \label{lem:dc-ch9-2-4-piecewise-assembly}
  \dcref{ch9:2.4}
  \lean{Riemannian.Variation.hasDerivAt_dcEnergy_eq_piecewise_first_variation, Riemannian.Variation.hasDerivAt_dcEnergy_eq_piecewise_first_variation_of_segment_formulas}
  \leanok
  \uses{def:dc-ch9-2-2-energy, lem:dc-ch9-2-4-segment}
  Let $0=t_0<\cdots<t_{k+1}=a$. Assume the energy density is integrable on every
  segment near $s_0$ and formula (1) has been established separately on each
  $[t_i,t_{i+1}]$. Then the energy on $[0,a]$ is the sum of the segment energies, its
  derivative is the sum of their derivatives, and the endpoint terms telescope to
  $$
  \left\langle V(a),\frac{dc}{dt}(a^-)\right\rangle
  -\left\langle V(0),\frac{dc}{dt}(0^+)\right\rangle
  -\sum_{i=1}^k\left\langle V(t_i),
    \frac{dc}{dt}(t_i^+)-\frac{dc}{dt}(t_i^-)\right\rangle.
  $$
  Thus the finite-additivity, differentiation, and jump-term algebra needed for the
  piecewise formula are formalized. The full proposition still has to extract the
  segment hypotheses and one-sided fields from an arbitrary variation.
\end{lemma}
\begin{proof}
  \leanok
  \uses{def:dc-ch9-2-2-energy, lem:dc-ch9-2-4-segment}
  Apply finite additivity of interval integrals, differentiate the resulting finite sum,
  and telescope consecutive endpoint terms. $\square$
\end{proof}

\begin{lemma}[proper smooth variations through a geodesic are stationary]
  \label{lem:dc-ch9-2-5-geodesic-energy}
  \dcref{ch9:2.5}
  \lean{Riemannian.Variation.hasDerivAt_dcEnergy_zero_of_geodesic, Riemannian.Variation.IsVariation.hasDerivAt_dcEnergy_zero_of_proper_geodesic, Riemannian.Variation.hasDerivAt_dcEnergy_zero_of_piecewise_geodesic_segment_formulas}
  \leanok
  \uses{lem:dc-ch9-2-4-segment, lem:dc-ch9-2-5-geodesic, def:dc-ch9-2-1}
  In the hypotheses of \cref{lem:dc-ch9-2-4-segment} on a segment $[a,b]$, suppose that
  the covariant derivative of the velocity of the base curve is zero and that the
  variational field vanishes at both endpoints. Then the energy of the variation is
  stationary at the base parameter:
  $$
  E'(s_0)=0.
  $$
\end{lemma}
\begin{proof}
  \leanok
  \uses{lem:dc-ch9-2-4-segment, lem:dc-ch9-2-5-geodesic}
  Formula (1) of \cref{lem:dc-ch9-2-4-segment} has three contributions: the integral
  against the covariant acceleration and the two endpoint pairings. The first is zero
  when the base curve is a geodesic, and the endpoint pairings are zero for a proper
  variation. Hence $E'(s_0)=0$.
\end{proof}

\begin{lemma}[telescoping the breakpoint terms]
  \label{lem:dc-ch9-2-8-jump-sum}
  \dcref{ch9:2.8}
  \lean{Riemannian.Variation.sum_segment_boundary_eq_endpoint_sub_sum_jump,
    Riemannian.Variation.sum_segment_boundary_eq_neg_sum_jump}
  \leanok
  Let $0=t_0<t_1<\cdots<t_{k+1}=a$ be a subdivision and write
  \[
    b_i^- = \left\langle V(t_i),\frac{DV}{dt}(t_i^-)\right\rangle,
    \qquad
    b_i^+ = \left\langle V(t_i),\frac{DV}{dt}(t_i^+)\right\rangle.
  \]
  Then the sum of the endpoint contributions from the $k+1$ segments is
  \[
    \sum_{i=0}^k (b_{i+1}^- - b_i^+)
      = b_{k+1}^- - b_0^+ - \sum_{i=1}^k (b_i^+ - b_i^-).
  \]
  In particular, for a proper variation the two outer terms vanish, and the
  segment boundaries give exactly the negative jump sum in formula (3).
\end{lemma}
\begin{proof}
  \leanok
  Induct on the number of internal subdivision points.  Adding the last
  segment replaces the old terminal term by the new terminal term and the
  difference $b_{k+1}^- - b_{k+1}^+$, which is the next jump contribution.
  The proper case follows by setting the two outer pairings equal to zero.
\end{proof}

\begin{lemma}[piecewise second-variation assembly under segment hypotheses]
  \label{lem:dc-ch9-2-8-piecewise-assembly}
  \dcref{ch9:2.8}
  \lean{Riemannian.Variation.hasDerivAt_deriv_dcEnergy_eq_piecewise_second_variation,
    Riemannian.Variation.hasDerivAt_deriv_dcEnergy_eq_piecewise_second_variation_of_proper,
    Riemannian.Variation.deriv_deriv_dcEnergy_eq_piecewise_second_variation_of_proper}
  \leanok
  \uses{lem:dc-ch9-2-8-segment, lem:dc-ch9-2-8-jump-sum, def:dc-ch9-2-2-energy}
  Fix a finite subdivision $\tau_0,\ldots,\tau_{k+1}$. Assume that, near the base
  parameter $s_0$, each segment energy is integrable and differentiable, and assume that
  the segment second-variation derivative is
  $$
  \frac{d}{ds}\,E_i'(s_0)=2\big((b^-_{i+1}-b^+_i)-B_i\big).
  $$
  Then finite additivity of energy and differentiation of a finite sum give the global
  derivative with endpoint terms and the internal jump sum:
  $$
  E''(s_0)=2\left(b^-_{k+1}-b^+_0-
    \sum_{i=1}^{k}(b^+_i-b^-_i)-\sum_{i=0}^{k}B_i\right).
  $$
  If the two outer pairings vanish, $b^+_0=b^-_{k+1}=0$ (as properness supplies
  in do Carmo's setting), this yields formula (3) with the subdivision jump term.
  The hypotheses are intentionally stated at the segment level:
  this lemma supplies the finite-subdivision assembly, while \cref{prop:dc-ch9-2-8} still
  needs the geometric/regularity bridge from an arbitrary variation to those hypotheses.
\end{lemma}
\begin{proof}
  \leanok
  \uses{lem:dc-ch9-2-8-segment, lem:dc-ch9-2-8-jump-sum}
  Apply finite additivity of the energy integral on the subdivision, differentiate the finite
  sum, and substitute the assumed segment identities. Reindexing the endpoint terms gives
  the displayed telescoping expression; properness removes the two outer pairings. $\square$
\end{proof}

\begin{lemma}[sectional Bonnet--Myers under explicit variation data]
  \label{lem:dc-ch9-3-3-sectional-analytic}
  \dcref{ch9:3.3}
  \lean{Riemannian.Variation.bonnetMyers_diameterBound_of_sectional_of_analytic}
  \leanok
  \uses{lem:dc-ch9-3-3-ricci-sectional, lem:dc-ch9-3-1-velocity-frame,
    lem:dc-ch9-3-1-index-sum-neg, thm:dc-ch7-2-8}
  Let $M$ be a connected, complete Riemannian manifold of dimension at least two, and let
  $r>0$. Assume that every sectional curvature is at least $1/r^2$. Suppose in addition that
  along every distance-realizing geodesic of length $\ell>\pi r$, the velocity-seeded parallel
  orthonormal frame and the sine fields $V_j(t)=\sin(\pi t)e_j(t)$ satisfy the explicit
  nonnegativity conclusion of the second-variation argument (the analytic variation data).
  Then $\operatorname{diam}(M)\le\pi r$ and $M$ is compact.

  This is the sectional-curvature specialization of the Bonnet--Myers assembly with its
  remaining variation regularity exposed; it makes no assertion about $\pi_1(M)$.
\end{lemma}
\begin{proof}
  \leanok
  By \cref{lem:dc-ch9-3-3-ricci-sectional}, the sectional lower bound gives the unnormalized
  Ricci lower bound $(n-1)/r^2$ in every orthonormal frame. If two points were farther apart
  than $\pi r$, Hopf--Rinow (\cref{thm:dc-ch7-2-8}) would provide a distance-realizing geodesic
  of length $\ell>\pi r$. The velocity-seeded frame and the sine index-form computation
  (\cref{lem:dc-ch9-3-1-velocity-frame,lem:dc-ch9-3-1-index-sum-neg}) would make the sum of
  the perpendicular index forms strictly negative, while the assumed analytic variation data
  make each summand nonnegative, a contradiction. Thus all distances are at most $\pi r$.
  Completeness and Hopf--Rinow then make the bounded manifold compact.
\end{proof}

\begin{remark}[Remark 2.11]
  \label{rem:dc-ch9-2-11}
  \dcref{ch9:2.11}
  It is possible to establish formulas for the first and second variations of the
  arc length function $L(s)$. The expressions are entirely analogous to those for the
  energy and, since they will not be used here, are not treated (see [dC 2]
  paragraph 5.4).
\end{remark}

\begin{remark}[Remark 2.6]
  \label{rem:dc-ch9-2-6}
  \dcref{ch9:2.6}
  \uses{prop:dc-ch9-2-5}
  \cref{prop:dc-ch9-2-5} furnishes a characterization of geodesics as critical points of
  the energy for all proper variations. In this sense geodesics can be thought of as
  solutions to a variational problem; such a characterization is not local but
  involves the behavior of the geodesic as a whole.
\end{remark}

\begin{remark}[Remark 3.5]
  \label{rem:dc-ch9-3-5}
  \dcref{ch9:3.5}
  In fact it is not necessary that $K$ be bounded away from zero but only that $K$
  not approach zero too fast. In this respect, see E. Calabi, "On Ricci curvatures
  and geodesics", Duke Math. J. 34 (1967), 667–676 and R. Schneider, "Konvexe
  Flächen mit langsam abnehmender Krümmung", Archiv der Math. 23 (1972), 650–654.
\end{remark}
